\chapter{谱不变量：应用}
\label{ch:spectral-applications}

尽管谱不变量的构造严重依赖伪全纯曲线的解析理论，尤其依赖辛流形的光滑结构，
但这些不变量本身却是 $C^0$ 型不变量：它们在 (OhM07) 所引入的 $C^0$-哈密顿拓扑
（亦见本书第 6.2 节）中连续，因而可以延拓到第 6.2 节所定义的拓扑哈密顿流。

本章说明谱不变量在辛拓扑研究中的若干应用。与霍费尔型的其他更直接的动力学不变量
相比，谱不变量的一项重要优势是其\emph{同伦不变性}；这使我们能够自然地把谱不变量推前到
哈密顿微分同胚群的万有覆盖空间，有时甚至还能降到群本身。恩托夫与波尔捷罗维奇
(EnP06) 对部分辛拟态的惊人构造突出体现了这一点：该构造只以谱不变量
$\rho(H;1)$ 的公理性质以及动力系统中取渐近平均这一自然运算为基础。他们在单调情形
下完成了这一构造，随后厄舍 (Ush10b) 将之推广到任意紧辛流形。

首先，我们将说明 (Oh05d) 中不变谱范数的构造，以及它在辛刚性问题和霍费尔几何中
测地线极小性问题上的应用。我们还将说明厄舍对波尔捷罗维奇以及拉隆德--麦克达夫
极小性猜想和最优能量--容量不等式的应用。其次，我们将自足地介绍恩托夫与
波尔捷罗维奇从谱不变量构造的部分辛拟态与拟同态，以及它们在辛相交问题中的应用。
最后，我们回到哈密顿同胚群的研究，并说明如何把所有这些构造延拓到第 6 章意义下的
连续哈密顿范畴。

看来，这些构造作为一个整体对辛拓扑的发展颇为重要；从纯辛拓扑的观点看，这一发展
正开始揭示格罗莫夫伪全纯曲线与弗洛尔同调之真正含义所笼罩的神秘面纱。

\section{哈密顿微分同胚的谱范数}
\label{sec:spectral-norm-hamiltonian-diffeomorphisms}

本节依照 (Oh05d) 说明哈密顿微分同胚的一种不变范数的构造；它称为\emph{谱范数}。
这里需要谨慎使用对应于量子上同调类 $1\in QH^*(M)$ 的谱不变量 $\rho(H;1)$。

回忆 $C^\infty([0,1]\times M)$ 上有自然对合
\[
 H\longmapsto\overline H;\qquad
 \overline H(t,x)=-H(t,\phi_H^t(x)),
\]
它也等价于取逆的对合
\[
 \lambda\longmapsto\lambda^{-1};\qquad
 \mathcal P^{\mathrm{ham}}(\operatorname{Symp}(M,\omega),\mathrm{id})
 \longrightarrow
 \mathcal P^{\mathrm{ham}}(\operatorname{Symp}(M,\omega),\mathrm{id}).
\]
我们关于这个对合将 $\rho(\mathord\cdot;1)$ 对称化，并定义函数
\[
 \gamma:C^\infty([0,1]\times M)\longrightarrow\mathbb R
\]
为
\begin{equation}
 \gamma(H)=\rho(H;1)+\rho(\overline H;1)
 \label{eq:22-gamma-hamiltonian}
\end{equation}
（定义域为 $C^\infty([0,1]\times M)$）。显然，对任意 $H$ 都有
$\gamma(H)=\gamma(\overline H)$。一般三角不等式
\[
 \rho(H;a)+\rho(F;b)\geq\rho(H\#F;a\cdot b)
\]
限制到 $a=b=1$，再结合定理 21.3.7 的归一化公理 $\rho(\mathrm{id};1)=0$，便给出
\begin{equation}
 \gamma(H)=\rho(H;1)+\rho(\overline H;1)\geq\rho(0;1)=0.
 \label{eq:22-gamma-nonnegative}
\end{equation}
这里 $a\cdot b$ 是量子上同调类 $a,b\in QH^*(M)$ 的量子乘积，$0$ 是零函数。

下面把 $\gamma(H)$ 与 $L^{(1,\infty)}$-范数
$\lVert H\rVert=\int_0^1\operatorname{osc}(H_t)\,dt$ 比较。回忆定义
\begin{equation}
 E^-(H)=\int_0^1-\min H\,dt,
 \qquad
 E^+(H)=\int_0^1\max H\,dt,
 \label{eq:22-hofer-lengths}
\end{equation}
\stfhSubIndex{I109}{Huofeierchangdu}{霍费尔长度（Hofer length）}{Fu}{负霍费尔长度（negative Hofer length）}
\stfhSubIndex{I110}{Huofeierchangdu}{霍费尔长度（Hofer length）}{Zheng}{正霍费尔长度（positive Hofer length）}
分别称为负霍费尔长度与正霍费尔长度。显然
$E^+(\overline H)=E^-(H)$。我们证明下面的引理。

\begin{lemma}
\label{lem:22-spectral-hofer-bounds}
对任意 $H$ 和 $0\neq a\in QH^*(M)$，有
\begin{equation}
 -E^+(H)+v(a)\leq\rho(H;a)\leq E^-(H)+v(a),
 \label{eq:22-spectral-valuation-bounds}
\end{equation}
其中 $v(a)$ 是 $a$ 的\termValuation{}。特别地，对任意经典上同调类
$b\in H^*(M)\hookrightarrow QH^*(M)$，任意哈密顿量 $H$ 都满足
\begin{equation}
 -E^+(H)\leq\rho(H;b)\leq E^-(H).
 \label{eq:22-classical-hofer-bounds}
\end{equation}
\end{lemma}

\begin{proof}
从不等式 (21.3.22) 出发，即
\[
 \int_0^1-\max(H-K)\,dt
 \leq\rho(H;a)-\rho(K;a)
 \leq\int_0^1-\min(H-K)\,dt.
\]
可将它改写成
\[
 \rho(K;a)+\int_0^1-\max(H-K)\,dt
 \leq\rho(H;a)\leq
 \rho(K;a)+\int_0^1-\min(H-K)\,dt.
\]
令 $K\to0$，得到
\begin{equation}
 \rho(0;a)+\int_0^1-\max H\,dt
 \leq\rho(H;a)\leq
 \rho(0;a)+\int_0^1-\min H\,dt.
 \label{eq:22-zero-hamiltonian-bound}
\end{equation}
由归一化公理，$\rho(0;a)=v(a)$，故 (21.3.22) 变为
\[
 v(a)-E^+(H)\leq\rho(H;a)\leq v(a)+E^-(H).
\]
经典上同调类 $b$ 满足 $v(b)=0$，于是由定义立刻得到
\eqref{eq:22-classical-hofer-bounds}。
\end{proof}

把~\eqref{eq:22-classical-hofer-bounds} 的右侧分别用于 $b=1$ 和
$H,\overline H$，便有 $\rho(H;1)\leq E^-(H)$ 以及
$\rho(\overline H;1)\leq E^-(\overline H)$。另一方面，对任意 $H$，
$E^-(\overline H)=E^+(H)$，所以
\begin{equation}
 \gamma(H)\leq\lVert H\rVert.
 \label{eq:22-gamma-hofer-bound}
\end{equation}
非负性~\eqref{eq:22-gamma-nonnegative} 引出如下定义。

\begin{definition}
\label{def:22-spectral-norm}
定义 $\gamma:\operatorname{Ham}(M,\omega)\to\mathbb R_+$ 为
\[
 \gamma(\phi):=\inf_{H\mapsto\phi}\gamma(H)
 =\inf_{H\mapsto\phi}\bigl(\rho(H;1)+\rho(\overline H;1)\bigr).
\]
\stfhIndex{I175}{Pufanshu}{谱范数（spectral norm）}
称 $\gamma(\phi)$ 为 $\phi$ 的\emph{谱范数}。
\end{definition}

下面的定理概括函数 $\gamma$ 的基本性质。在本小节余下部分，我们将证明这些陈述；
其中最非平凡的非退化性留到下一小节。

\begin{theorem}
\label{thm:22-spectral-norm-properties}
设 $\gamma$ 如上。则
$\gamma:\operatorname{Ham}(M,\omega)\to\mathbb R_+$ 定义一个不变范数，即它满足：
\begin{enumerate}
 \item （非退化性）$\phi=\mathrm{id}$ 当且仅当 $\gamma(\phi)=0$；
 \item （辛不变性）对任意辛微分同胚 $\eta$，
 $\gamma(\eta^{-1}\phi\eta)=\gamma(\phi)$；
 \item （三角不等式）$\gamma(\phi\psi)\leq\gamma(\phi)+\gamma(\psi)$；
 \item （取逆不变性）$\gamma(\phi^{-1})=\gamma(\phi)$；
 \item $\gamma(\phi)\leq\lVert\phi\rVert$。
\end{enumerate}
\end{theorem}

\begin{proof}[辛不变性的证明]
回忆谱不变量的辛不变性
$\rho(H;a)=\rho(\eta_*H;\eta_*a)$。把它用于 $a=1$，并利用
$\eta_*1=1$，得到
\[
 \begin{aligned}
 \gamma(\phi)
 &=\inf_{H\mapsto\phi}\bigl(\rho(H;1)+\rho(\overline H;1)\bigr)\\
 &=\inf_{H\mapsto\phi}
 \bigl(\rho(\eta_*H;1)+\rho(\overline{\eta_*H};1)\bigr)
 =\gamma(\eta^{-1}\phi\eta).
 \end{aligned}
\]
\end{proof}

\begin{proof}[三角不等式的证明]
先回忆三角不等式
\begin{equation}
 \rho(H\#K;1)\leq\rho(H;1)+\rho(K;1)
 \label{eq:22-triangle-forward}
\end{equation}
以及
\begin{equation}
 \rho(\overline K\#\overline H;1)
 \leq\rho(\overline K;1)+\rho(\overline H;1).
 \label{eq:22-triangle-inverse}
\end{equation}
将~\eqref{eq:22-triangle-forward} 与~\eqref{eq:22-triangle-inverse} 相加，得到
\begin{align}
 \rho(H\#K;1)+\rho(\overline{H\#K};1)
 &=\rho(H\#K;1)+\rho(\overline K\#\overline H;1)\notag\\
 &\leq\bigl(\rho(H;1)+\rho(\overline H;1)\bigr)
 +\bigl(\rho(K;1)+\rho(\overline K;1)\bigr).
 \label{eq:22-gamma-triangle-hamiltonians}
\end{align}
现令 $H\mapsto\phi$、$K\mapsto\psi$。因为 $H\#K$ 生成 $\phi\psi$，所以
\[
 \gamma(\phi\psi)\leq\gamma(H\#K)
 =\rho(H\#K;1)+\rho(\overline{H\#K};1),
\]
进而由~\eqref{eq:22-gamma-triangle-hamiltonians} 得
\[
 \gamma(\phi\psi)\leq
 \bigl(\rho(H;1)+\rho(\overline H;1)\bigr)
 +\bigl(\rho(K;1)+\rho(\overline K;1)\bigr).
\]
让右端对所有 $H\mapsto\phi$ 与 $K\mapsto\psi$ 取下确界，即证性质 (3)。
\end{proof}

\begin{proof}[取逆不变性的证明]
由 $\gamma$ 的定义关于映射 $\phi\mapsto\phi^{-1}$ 对称，结论立即成立。
\end{proof}

\begin{proof}[性质 (5) 的证明]
令~\eqref{eq:22-gamma-hofer-bound} 对 $H\mapsto\phi$ 取下确界，便得
$\gamma(\phi)\leq\lVert\phi\rVert$。
\end{proof}

现在只剩证明 $\gamma$ 的非退化性。(Oh05d) 中的原证明使用了特定伪全纯曲线的一条
相当精细的存在性定理。尽管这条存在性定理本身对刚性问题有重要应用，我们仍将把非退化性
作为一个更熟悉的能量--容量型不等式的推论来给出另一证明；这个证明由厄舍 (Ush10a)
提出。它较为间接，需要暂时转向研究 $\operatorname{Ham}(M,\omega)$ 上的霍费尔几何。
（麦克达夫与萨拉蒙在 (MSa04) 中还给出了另一证明。）

\begin{remark}
\label{rem:22-medium-hofer-norm}
事实上，如果把 $\rho(\mathord\cdot;1)$ 的同伦不变性纳入
$\gamma(\phi)\leq\lVert\phi\rVert$ 的证明，就可以如下改进该不等式。定义\emph{中等霍费尔}
范数
\[
 \lVert\phi\rVert_{\mathrm{med}}
 :=\inf\{\lVert h\rVert\mid\pi(h)=\phi\},
\]
其中对 $h\in\widetilde{\operatorname{Ham}}(M,\omega)$ 定义
\[
 \lVert h\rVert:=\inf\{\lVert H\rVert
 \mid h=[\phi_H]\in\widetilde{\operatorname{Ham}}(M,\omega)\}.
\]
于是上述证明给出 $\gamma(H)\leq\lVert[\phi_H]\rVert$，从而
\[
 \gamma(\phi)\leq\lVert\phi\rVert_{\mathrm{med}}.
\]
更多相关讨论见 (Oh09a) 第 11 节。
\end{remark}

\stfhIndex{I129}{Zhongdenghuofeierfanshu}{中等霍费尔范数（medium Hofer norm）}
\section{霍费尔测地线与周期轨道}
\label{sec:hofer-geodesics-periodic-orbits}

本节研究哈密顿路径在霍费尔距离下的长度极小性。我们先沿用\personBialy{}与\personPolterovich{}
(BP94) 以及拉隆德与麦克达夫 (LM95b) 的做法，引入一个必要判据，再介绍厄舍 (Ush10a)
对所谓极小性猜想的证明。同时，我们也将依照厄舍的论证证明
$\operatorname{Ham}(M,\omega)$ 上 $\gamma$-范数的非退化性。

\subsection{准自治哈密顿量}

$H$ 的 $L^{(1,\infty)}$-范数
$\lVert H\rVert=\int_0^1(\max H_t-\min H_t)\,dt$ 可等同于路径
$\phi_H:t\mapsto\phi_H^t$ 的芬斯勒长度
\[
 \operatorname{leng}(\phi_H)=\int_0^1
 \bigl(\max_xH(t,\phi_H^t(x))-\min_xH(t,\phi_H^t(x))\bigr)\,dt,
\]
其中
\[
 T_{\mathrm{id}}\operatorname{Ham}(M,\omega)
 \cong C^\infty(M)/\mathbb R
\]
上的相伴巴拿赫范数正是
\[
 \lVert h\rVert=\operatorname{osc}(h)=\max h-\min h.
\]
映射 $t\mapsto H(t,\phi_H^t(\mathord\cdot))$ 可视为
$\operatorname{Ham}(M,\omega)$ 上路径 $t\mapsto\phi_H^t$ 的切向量。

\begin{definition}
\label{def:22-hofer-path-metric}
考虑定义在
$\mathcal P^{\mathrm{ham}}(\operatorname{Symp}(M,\omega),\mathrm{id})$ 上的度量
$d:\mathcal P^{\mathrm{ham}}(\operatorname{Symp}(M,\omega),\mathrm{id})
\to\mathbb R_+$：
\[
 d(\lambda,\mu):=\operatorname{leng}(\lambda^{-1}\circ\mu),
\]
其中 $\lambda^{-1}\circ\mu$ 是哈密顿路径
$t\in[0,1]\mapsto\lambda(t)^{-1}\mu(t)$。路径空间
$\mathcal P^{\mathrm{ham}}(\operatorname{Symp}(M,\omega),\mathrm{id})$
上的霍费尔拓扑即由此度量诱导。
\end{definition}

\begin{exercise}
\label{ex:22-hofer-quotient-topology}
考虑评价映射
\[
 \operatorname{ev}_1:
 \mathcal P^{\mathrm{ham}}(\operatorname{Symp}(M,\omega),\mathrm{id})
 \longrightarrow\operatorname{Ham}(M,\omega).
\]
假设 $M$ 紧。证明：上述 $\operatorname{Ham}(M,\omega)$ 上的霍费尔拓扑，是使评价
映射 $\lambda\mapsto\lambda(1)$ 连续的最强拓扑。
\end{exercise}

该拓扑的非平凡性来自霍费尔范数函数
$\lVert\mathord\cdot\rVert:\operatorname{Ham}(M,\omega)\to\mathbb R_+$
的非退化性。霍费尔 (H93) 证明：只要相应哈密顿方程不存在\emph{周期小于或等于一}的
非常值周期轨道，$\mathbb C^n$ 上任意紧支\emph{自治}哈密顿量的路径便是长度极小的。
(En00), (MSl01) 与 (Oh02)--(Oh05b) 在附加假设每个不动点处的线性化流不是过扭的
条件下推广了这个结果；亦即，该线性化流没有周期小于一的闭轨道。

\personBialy{}--\personPolterovich{}在 (BP94)、\personLalonde{}--\personMcDuff{}在 (LM95b) 中证明：任意长度极小
（分别地，局部长度极小）的哈密顿路径，都由\emph{准自治}（分别地，\emph{局部准自治}）哈密顿路径
生成。

\begin{definition}
\label{def:22-quasi-autonomous}
若存在两个在时间 $t\in[0,1]$ 内保持不动的点 $x_-,x_+\in M$，使得对所有
$t\in[0,1]$ 都有
\[
 H(t,x_-)=\min_xH(t,x),\qquad
 H(t,x_+)=\max_xH(t,x),
\]
\stfhSubIndex{I106}{Hamidunliang}{哈密顿量（Hamiltonian）}{Zhunzizhi}{准自治（quasi-autonomous）}
则称哈密顿量 $H$ 为\emph{准自治的}。
\end{definition}

下述定义一直用于研究哈密顿路径的长度极小性质。

\begin{definition}
\label{def:22-twisted-periodic-point}
设 $H:[0,1]\times M\to\mathbb R$ 是不必时间周期的哈密顿量，$\phi_H^t$ 是其哈密顿流。
\begin{enumerate}
 \item 若 $\phi_H^T(p)=p$，则称点 $p\in M$ 是\emph{时间-$T$ 周期点}。若轨道
 $t\in[0,T]\mapsto\phi_H^t(p)$ 可缩，则称之为\emph{可缩时间-$T$ 周期轨道}。
 \item 若 $H$ 在 $t\in[0,T]$ 内有不变临界点 $p$，且其线性化流
 \[
  d\phi_H^t(p):T_pM\longrightarrow T_pM,\qquad t\in[0,T],
 \]
 有周期小于或等于 $T$ 的闭轨道，则称 $p$ 作为时间-$T$ 周期轨道是\emph{过扭的}；
 否则称为\emph{欠扭的}。若线性化流此外只有原点作为不动点，则称该不动点为
 \emph{通有欠扭的}。
\end{enumerate}
\end{definition}

注意，通有欠扭性是哈密顿量 $H$ 的 $C^2$-稳定性质。

为引出后面的讨论，我们回忆\personUstilovsky{}--\personLalonde{}--\personMcDuff{}关于测地线稳定性的必要
条件，但略去证明。

\begin{theorem}[\personUstilovsky{} (Ust96)，\personLalonde{}--\personMcDuff{} (LM95b)]
\label{thm:22-stable-geodesic-fixed-points}
设 $\phi$ 是通有的，即其所有不动点均孤立。则从恒等映射到 $\phi$ 的任意稳定测地线
$\phi^t$（$0\leq t\leq1$），必有至少两个欠扭不动点。
\end{theorem}

下面的猜想由波尔捷罗维奇提出（(Po01) 猜想 12.6.D），有时称为极小性猜想。
（亦见 (Po01), (LM95b), (MSl01)。）

\begin{conjecture}[极小性猜想]
\label{conj:22-minimality}
任意不存在周期小于或等于一的可缩周期轨道的自治哈密顿路径，都在其端点固定的路径同伦类
中是霍费尔长度极小的。
\end{conjecture}

恩托夫 (En00) 的下述定义简化了术语。

\begin{definition}
\label{def:22-slow-hamiltonian}
\stfhSubIndex{I107}{Hamidunliang}{哈密顿量（Hamiltonian）}{Man}{慢（slow）}
若哈密顿量 $H$ 没有周期小于或等于一的可缩周期轨道，则称 $H$ 为\emph{慢的}。
\end{definition}

在相伴自治哈密顿路径的不动点线性化流满足各种条件时，多位作者证明了极小性猜想。
厄舍 (Ush10a) 在不对线性化流附加任何限制的情况下完整证明了此猜想。他使用
(Oh05b) 中的长度极小判据以及哈密顿量的一条逼近定理。（相关结果亦见 (Sch06)。）
下面两个小节将介绍厄舍的证明。

\subsection{\texorpdfstring{通过 $\rho(H;1)$ 的长度极小判据}{通过 rho(H;1) 的长度极小判据}}

我们用 (Oh05b) 中的谱不变量 $\rho(H;1)$ 描述哈密顿路径长度极小性质的一个简单判据。
该判据类似于 (H93) 以及 $\mathbb C^n$ 情形的 (BP94) 中所用的判据。

为描述这一判据，把~\eqref{eq:22-hofer-lengths} 中定义的范数分解为
\[
 \lVert H\rVert=E^-(H)+E^+(H)
\]
会很有用。

\begin{proposition}
\label{prop:22-negative-length-minimizing}
设 $G:[0,1]\times M\to\mathbb R$ 是满足
\begin{equation}
 \rho(G;1)=E^-(G)
 \label{eq:22-negative-length-criterion}
\end{equation}
的任意哈密顿量。则 $G$ 在其端点固定的同伦类中使负霍费尔长度极小。特别地，若此外
$\rho(\overline G)=E^-(\overline G)=E^+(G)$，则 $G$ 必为准自治的。
\footnote{译者注：原书命题首项结论印作“Then $H$”，但命题只给定了 $G$；译文据证明
中的 $F\sim G$ 改作 $G$。}
\end{proposition}

\begin{proof}
设 $F$ 是满足 $F\sim G$ 的任意哈密顿量。由~\eqref{eq:22-negative-length-criterion}、
(21.6.39)（取 $a=1$）以及~\eqref{eq:22-classical-hofer-bounds}，依次得到
\[
 E^-(G)=\rho(G;1)=\rho(F;1)\leq E^-(F).
\]
这证明了第一项陈述。第二项陈述于是由\personBialy{}--\personPolterovich{}、\personUstilovsky{}以及
拉隆德--麦克达夫的极小性判据得出。
\end{proof}

另一方面，\emph{若 $G$ 是一周期的}，就可以考虑相伴作用量泛函 $\mathcal A_G$。对于准自治
周期哈密顿量 $G$，$\mathcal A_G$ 有两个显然的临界值：
\[
 \mathcal A_G([x_-,\widetilde x_-])
 =\int_0^1-G(t,x_-)\,dt,
 \qquad
 \mathcal A_G([x_+,\widetilde x_+])
 =\int_0^1-G(t,x_+)\,dt.
\]
要点在于，这些作用量值分别与 $\lVert H\rVert$ 的负、正部分
\[
 E^-(G)=\int_0^1-\min G_t\,dt,
 \qquad
 E^+(G)=\int_0^1\max G_t\,dt
\]
一致。


\subsection{基本弗洛尔闭链的构造}

现在引入同调本质性的一个基本概念，即\emph{在柯西--黎曼流下的不可下推性质}。
这是\emph{链层弗洛尔理论}的核心，特别是谱不变量的构造 (Oh05c) 及其应用
(Oh05b, Oh05d, Oh09a) 的核心。这个概念最早在 (Oh97b) 关于余切丛中正合拉格朗日
子流形的语境下隐含地引入弗洛尔理论；随后在一般哈密顿微分同胚的语境下，作为
紧弗洛尔闭链的概念明确出现于 (Oh02, Oh09a)。

由定义 21.5.2 中紧弗洛尔闭链的定义，若量子同调类满足 $[\alpha]=a^\flat$，则
$(H,J)$ 的任意紧弗洛尔闭链都实现数值 $\rho(H;a)$，即
$\lambda_{(H,J)}(\alpha)=\rho(H;a)$。证明这种闭链存在并非易事。下述定理在 (Oh09a)
中由几何连续性方法证明，在 (Ush08) 中由纯代数方法证明。

\begin{theorem}[吴龙根 (Oh09a)，厄舍 (Ush08)]
\label{thm:22-tight-floer-cycle-exists}
对任意 $0\neq a\in QH^*(M)$，都存在一个与 $a$ 相伴的紧弗洛尔闭链。
\end{theorem}

此定理的一般证明相当微妙（除非 $(M,\omega)$ 是有理的），超出本书范围。因此请读者
参阅原始论文。

与命题~\ref{prop:22-negative-length-minimizing} 中的长度极小判据有关，对满足慢性条件的
准自治哈密顿量 $G$，关键在于构造一个与 $a=1$ 相伴、层级
$\lambda_{(G,J)}(\alpha)=E^-(G)$ 的紧弗洛尔闭链。

然而，一般哈密顿量的弗洛尔闭链定义因其表述方式而相当抽象复杂，故通常不易用于实际
应用。另一方面，给定 $(M,\omega)$ 的基本闭链所对应的弗洛尔闭链，可以通过弗洛尔链
映射传递莫尔斯基本闭链来系统研究。

\begin{definition}
\label{def:22-fundamental-floer-cycle}
若弗洛尔闭链 $\alpha$ 的同调类满足 $[\alpha]=1^\flat$，则称 $\alpha$ 为\emph{基本闭链}。
\end{definition}

类似地，莫尔斯基本闭链在莫尔斯理论中定义。当哈密顿路径满足准自治性质之类的特定
性质时，我们希望在研究这些传递所得的闭链时利用这一几何性质。需要强调的是，一般而言，
\emph{这个过程只在弗洛尔链层面而非临界点层面良定义}。这是因为在弗洛尔链映射下，一个临界点
被传递为一条弗洛尔链，而不仅是另一个临界点。

选择莫尔斯函数 $f$，并取\emph{充分小的} $\epsilon>0$ \emph{使定理 19.8.1 成立}；再取
$-\epsilon f$ 的莫尔斯基本闭链 $\alpha$。于是也可把 $\alpha$ 视为
$\epsilon f$ 的弗洛尔闭链。通过弗洛尔链映射 $h_{\mathcal L}$ 传递 $\alpha$，定义
$H$ 的基本弗洛尔闭链
\[
 \alpha_H:=h_{\mathcal L}(\alpha)\in CF(H),
\]
其中 $h_{\mathcal L}$ 是任意连接 $\epsilon f$ 与 $H$ 的给定同伦 $\mathcal L$ 上的
弗洛尔链映射。关于极小性猜想，典范给出的线性路径
\[
 \mathcal L:s\longmapsto(1-s)\epsilon f+sH
\]
起着重要作用。还要注意，这条闭链依赖于莫尔斯函数 $f$ 的选择。一般而言，即使 $H$
是准自治的，也不能期待这条闭链是紧的。

\stfhIndex{I116}{Keerman-Lalongde}{克尔曼--拉隆德（Kerman--Lalonde）}
在这方面，克尔曼--拉隆德 (KL03) 证明：在下面两项要求下，
\begin{enumerate}
 \item 准自治哈密顿量 $H$ 有唯一的非退化整体极小点 $x_-$，且该点对所有
 $t\in[0,1]$ 均欠扭；
 \item 出现在上述线性路径 $\mathcal L$ 中的莫尔斯函数 $f$
 \begin{itemize}
  \item 以 $x_-$ 为整体极小点；
  \item 且对所有其他临界点 $x_j$ 满足
  \begin{equation}
   f(x_-)=0,\qquad f(x_-)<f(x_j),
   \label{eq:22-morse-minimum-condition}
  \end{equation}
 \end{itemize}
\end{enumerate}
\stfhSubIndex{I085}{Fuluoerbilian}{弗洛尔闭链（Floer cycle）}{Jibenbilian}{基本闭链（fundamental cycle）}
从 $f$ 的莫尔斯基本闭链经 $\mathcal L$ 传递得到的相应基本弗洛尔闭链具有如下良好
性质。该结果在 (KL03) 中对辛非球面流形证明，在 (Oh05b) 中推广到一般辛流形。

\begin{proposition}
\label{prop:22-canonical-fundamental-cycle}
设 $H$ 是通有的一周期哈密顿量，使得 $H_t$ 在 $x_-$ 处有唯一的非退化整体极小值，
而 $x_-$ 对所有 $t\in[0,1]$ 都是不动且欠扭的。设 $f:M\to\mathbb R$ 是以 $x_-$
为唯一整体极小点、并满足 $f(x_-)=0$ 的莫尔斯函数。则典范基本闭链具有表达式
\begin{equation}
 \alpha_H=[x_-,\widetilde x_-]+\beta\in CF(H),
 \label{eq:22-fundamental-cycle-expression}
\end{equation}
其中某条弗洛尔链 $\beta\in CF(H)$ 的层级满足
\begin{equation}
 \lambda_H(\beta)<\lambda_H([x_-,\widetilde x_-])
 =\int_0^1-H(t,x_-)\,dt.
 \label{eq:22-beta-level}
\end{equation}
特别地，$\alpha_H$ 的层级满足
\begin{align}
 \lambda_H(\alpha_H)
 &=\lambda_H([x_-,\widetilde x_-])
 \label{eq:22-fundamental-cycle-level}\\
 &=\int_0^1-H(t,x_-)\,dt
 =\int_0^1-\min H_t\,dt.\notag
\end{align}
\end{proposition}

\begin{proof}
证明基于下述引理（见 (KL03) 命题 5.2）。提醒读者，我们的符号约定与 (KL03) 不同。

\begin{lemma}
\label{lem:22-kerman-lalonde}
设 $H,f$ 如上。则对所有充分小的 $\epsilon>0$，定义为
\[
 G^H(t,x)=H(t,x_-)+\epsilon f(x)
\]
的函数 $G^H$ 满足
\[
 G^H(t,x_-)=H(t,x_-),\qquad G^H(t,x)\leq H(t,x)
\]
对所有 $(t,x)$ 成立，且等号只在 $x_-$ 处成立。
\end{lemma}

\begin{exercise}
\label{ex:22-prove-kerman-lalonde}
证明这个引理。
\end{exercise}

现在回到命题的证明。因为 $x_-$ 是 $H$ 与 $f$ 的欠扭不动极小点，所以有
康利--曾德尔指标
\[
 \mu_H([x_-,w_{x_-}])=\mu_{\epsilon f}([x_-,w_{x_-}])(=-n),
\]
从而模空间
$\mathcal M_{\mathcal L}([x_-,w_{x_-}],[x_-,w_{x_-}])$ 为零维。由于常值映射
$u(\tau,t)\equiv x_-$ 总在其中，故
\[
 \mathcal M_{\mathcal L}([x_-,w_{x_-}],[x_-,w_{x_-}])\neq\varnothing.
\]

任取
$u\in\mathcal M_{\mathcal L}([x_-,w_{x_-}],[x_-,w_{x_-}])$。注意，对线性同伦
\[
 \mathcal L:s\longmapsto(1-s)\epsilon f+sH
\]
的弗洛尔连续性方程，即使把同伦替换为
\[
 \mathcal L':s\longmapsto(1-s)G^H+sH,
\]
也保持不变。这是因为 $G^H$ 在 $\epsilon f$ 上增加的项 $H(t,x_-)$ 不依赖于
$x\in M$，因而
\[
 X_{\epsilon f}\equiv X_{G^H}.
\]
所以 $u$ 也是同伦 $\mathcal L'$ 下连续性方程的解。由此得到恒等式
\footnote{译者注：原书式 (22.2.16) 的第二个积分在括号内、外各印了一组
$dt\,d\tau$；译文按连续性能量恒等式只保留括号外的一组微分。}
\begin{align}
 \int\left\lVert\frac{\partial u}{\partial\tau}\right\rVert_{J^\rho}^2
 \,dt\,d\tau
 &=\mathcal A_{G^H}([x_-,w_{x_-}])
 -\mathcal A_H([x_-,w_{x_-}])\notag\\
 &\quad-\int_{-\infty}^{\infty}\rho'(\tau)
 \left(\int_0^1\bigl(H(t,u(\tau,t))-G^H(t,u(\tau,t))\bigr)\,dt\right)d\tau.
 \label{eq:22-continuation-energy-identity}
\end{align}
因为
\[
 \mathcal A_H([x_-,w_{x_-}])
 =\mathcal A_{G^H}([x_-,w_{x_-}])
 =\int_0^1-\min H\,dt
\]
且 $G^H\leq H$，故~\eqref{eq:22-continuation-energy-identity} 右端非正。因此
$\mathcal M_{\mathcal L}([x_-,w_{x_-}],[x_-,w_{x_-}])$ 只含常值解 $u\equiv x_-$。
特别地，这给出 $h_{\mathcal L}$ 的矩阵系数
\[
 \langle[x_-,w_{x_-}],h_{\mathcal L}([x_-,w_{x_-}])\rangle
 =\#\mathcal M_{\mathcal L}([x_-,w_{x_-}],[x_-,w_{x_-}])=1.
\]

现考虑 $\alpha_H$ 的任意其他生成元
\begin{equation}
 [z,w]\in\operatorname{supp}\alpha_H,
 \qquad [z,w]\neq[x_-,w_{x_-}].
 \label{eq:22-other-generator}
\end{equation}
由 $h_{\mathcal L}$ 与 $\alpha_H=h_{\mathcal L}(\alpha)$ 的定义，存在生成元
$[x,w_x]\in\operatorname{supp}\alpha$ 使得
\[
 \mathcal M_{\mathcal L}([x,w_x],[z,w])\neq\varnothing.
\]
于是对任意 $u\in\mathcal M_{\mathcal L}([x,w_x],[z,w])$，有恒等式
\[
 \begin{aligned}
 \mathcal A_H([z,w])-\mathcal A_{G^H}([x,w_x])
 &=-\int\left\lVert\frac{\partial u}{\partial\tau}\right\rVert_{J^\rho}^2
 \,dt\,d\tau\\
 &\quad-\int_{-\infty}^{\infty}\rho'(\tau)
 \left(\int_0^1(H(t,u(\tau,t))-G^H(t,u(\tau,t)))\,dt\right)d\tau.
 \end{aligned}
\]
由于 $-\lVert\partial u/\partial\tau\rVert_{J^\rho}^2\leq0$ 且 $G^H\leq H$，有
\begin{equation}
 \mathcal A_H([z,w])\leq\mathcal A_{G^H}([x,w_x]),
 \label{eq:22-action-comparison}
\end{equation}
等号只在 $u$ 为定常解时成立。分两种情况：$x=x_-$，以及
$[x_j,w_{x_j}]\in\operatorname{supp}\alpha$ 中 $x=x_j\neq x_-$。

第一种情形下，\emph{由假设} $[z,w]\neq[x_-,w_{x_-}]$，$u$ 不可能为常值，故
\eqref{eq:22-action-comparison} 中为严格不等式：
\begin{equation}
 \mathcal A_H([z,w])<\mathcal A_{G^H}([x_-,w_{x_-}]).
 \label{eq:22-strict-action-comparison}
\end{equation}
第二种情形下，某个 $x_j\neq x_-$ 满足~\eqref{eq:22-action-comparison}。注意，
\eqref{eq:22-morse-minimum-condition} 中的不等式 $f(x_-)<f(x_j)$ 等价于
\[
 \mathcal A_{G^H}([x_j,w_{x_j}])
 <\mathcal A_{G^H}([x_-,w_{x_-}]).
\]
它与~\eqref{eq:22-action-comparison} 一起再次给出
\eqref{eq:22-strict-action-comparison}。另一方面，因为
$G^H(t,x_-)=H(t,x_-)$，还有
\[
 \mathcal A_{G^H}([x_-,w_{x_-}])=\mathcal A_H([x_-,w_{x_-}]).
\]
综上，对任意
$[z,w]\in\operatorname{supp}\alpha_H$ 且 $[z,w]\neq[x_-,w_{x_-}]$，都有
\[
 \mathcal A_H([z,w])<\mathcal A_H([x_-,w_{x_-}])
 =\int_0^1-H(t,x_-)\,dt.
\]
命题得证。
\end{proof}

\subsection{自治哈密顿量的情形}

为证明极小性猜想，按照命题~\ref{prop:22-negative-length-minimizing} 给出的判据，只需证明
$\mathcal A_G([x_-,\widetilde x_-])=E^-(G)$ 与极小极大值 $\rho(G;1)$ 相等
（并类似证明 $\mathcal A_{\overline G}([x_+,\widetilde x_+])=E^-(\overline G)$
与 $\rho(\overline G;1)$ 相等）。后一事实是下述定理的直接推论。

\begin{theorem}
\label{thm:22-slow-fundamental-cycle-tight}
设 $G$ 是任意慢自治哈密顿量。则命题~\ref{prop:22-canonical-fundamental-cycle} 中构造的
典范基本闭链 $\alpha_G$ 是紧的，即
\[
 \rho(G;1)=\lambda_G(\alpha_G).
\]
特别地，$\rho(G;1)=E^-(G)$。类似地，
$\rho(\overline G;1)=E^-(\overline G)=E^+(G)$。
\end{theorem}

本节将证明这个定理。在开始证明前，先对其历史作一些说明。在以往多篇文献中，人们在
$G$ 满足下述条件或其变体时研究极小性猜想：
\begin{enumerate}
 \item $G$ 是慢的；
 \item 它有通有欠扭的极大点和极小点；
 \item 它的所有临界点在弗洛尔理论意义下非退化（即 $X_G$ 在每个临界点处的线性化流
 只有零点作为周期轨道）。
\end{enumerate}
例如参见 (MSl01), (En00), (Oh02), (KL03), (Oh05b)。条件 (2) 与 (3) 的出现，
其实是研究方法所造成的人为现象：该方法要使用某种形式的伪全纯曲线模空间，因而必须研究
模空间的横截性问题。通过这类研究可得的最强结果概括如下。

\begin{theorem}[定理 (Oh05b)]
\label{thm:22-oh-minimality}
设 $(M,\omega)$ 是任意闭辛流形，$G$ 是满足下列条件的自治哈密顿量：
\begin{enumerate}
 \item 它没有“周期为一”的非常值可缩周期轨道；
 \item 它有通有欠扭的极大点和极小点；
 \item 它的所有临界点在弗洛尔理论意义下非退化（即 $X_G$ 在每个临界点处的线性化流
 只有零点作为周期轨道）。
\end{enumerate}
则单参数子群 $\phi_G^t$ 对 $0\leq t\leq1$ 在其端点固定的同伦类中是长度极小的。
\end{theorem}

这个定理又是下述结果的直接推论。

\begin{theorem}
\label{thm:22-tight-under-twisted-cycle}
设 $G$ 是满足定理~\ref{thm:22-oh-minimality} 所列条件的任意慢自治哈密顿量。则命题
~\ref{prop:22-canonical-fundamental-cycle} 中构造的典范基本闭链 $\alpha_G$ 是紧的，即
\[
 \rho(G;1)=\lambda_G(\alpha_G)
 \bigl(=-G(x_-)=E^-(G)\bigr).
\]
\end{theorem}

\begin{proof}
注意，定理~\ref{thm:22-oh-minimality} 的条件特别蕴含 $G$ 非退化。固定一个不依赖时间的
$G$-正则殆复结构 $J_0$。

设 $\alpha$ 与典范基本弗洛尔闭链 $\alpha_G$ 同调，即
\begin{equation}
 \alpha=\alpha_G+\partial_G(\gamma)
 \label{eq:22-homologous-fundamental-cycle}
\end{equation}
对某条弗洛尔链 $\gamma\in CF_*(G)$ 成立。当 $G$ 自治且 $J\equiv J_0$ 不依赖于
$t$ 时，不存在落在 $[x_-,\widetilde x_-]$ 的非定常且不依赖于 $t$ 的
$\mathcal A_G$ 流线；因为任何这样的流线都来自 $G$ 的负莫尔斯梯度流，而 $x_-$
是 $G$ 的极小点。因此，任何落在 $[x_-,\widetilde x_-]$ 的非定常弗洛尔流线 $u$
必依赖于 $t$。由于假设 $G$ 不存在周期为一的非常值可缩周期轨道，
$\mathcal A_G$ 的任一临界点都形如
\[
 [x,w],\qquad x\in\operatorname{Crit}G.
\]

设 $u$ 是从 $[x,w]$（$x\in\operatorname{Crit}G$）出发、满足
\begin{equation}
 \mu([x,w])-\mu([x_-,\widetilde x_-])=1
 \label{eq:22-index-one-trajectory}
\end{equation}
的流线，并以
$\widehat{\mathcal M}_{(G,J_0)}([x,w],[x_-,\widetilde x_-])$ 表示相应的连接流线
弗洛尔模空间。一般指标公式 (18.4.24) 给出
\begin{equation}
 \mu([x,w])=\mu([x,w_x])-2c_1([w]).
 \label{eq:22-index-capping-formula}
\end{equation}
分别考虑 $c_1([w])=0$ 与 $c_1([w])\neq0$ 两种情形。若
$c_1([w])\neq0$，由~\eqref{eq:22-index-one-trajectory} 与
~\eqref{eq:22-index-capping-formula} 得 $x\neq x_-$。这意味着任意这样的流线都带有
（局部）自由的 $S^1$-作用（见第 19.8 节），即模空间
\[
 \mathcal M_{(G,J_0)}([x,w],[x_-,\widetilde x_-])
 =\widehat{\mathcal M}_{(G,J_0)}([x,w],[x_-,\widetilde x_-])/\mathbb R
\]
及其稳定映射紧化
$\overline{\mathcal M}_{(G,J_0)}([x,w],[x_-,\widetilde x_-])$ 都带有\emph{无不动点}的局部
自由 $S^1$-作用。于是由 (FHS95) 的 $S^1$-等变横截性定理，对适当选择的自治
$J_0$，模空间
$\mathcal M_{(G,J_0)}([x,w],[x_-,\widetilde x_-])$ 为空。这是因为在假设
~\eqref{eq:22-index-one-trajectory} 下，其商空间的虚维数为 $-1$。关于这一
$S^1$-不变正则化过程的更多说明，参见 (FHS95)。

现在考虑 $c_1([w])=0$ 的情形。首先，
~\eqref{eq:22-index-one-trajectory} 与~\eqref{eq:22-index-capping-formula}
蕴含 $x\neq x_-$。另一方面，若 $x\neq x_-$，上述同一论证表明扰动后的模空间为空。

因此，不存在任何落在 $[x_-,\widetilde x_-]$ 的指标 $1$ 流线。所以不论如何选择
$\gamma$，~\eqref{eq:22-homologous-fundamental-cycle} 中的 $\partial_G(\gamma)$
都不能消去闭链
\[
 \alpha_G=[x_-,\widetilde x_-]+\partial_G(\gamma)
\]
中的项 $[x_-,\widetilde x_-]$，故由层级 $\lambda_G$ 的定义有
\[
 \lambda_G(\alpha)\geq\lambda_G([x_-,\widetilde x_-]).
\]
结合~\eqref{eq:22-fundamental-cycle-level} 即完成证明。
\end{proof}

为把定理~\ref{thm:22-oh-minimality} 推广到极小性猜想中的一般慢哈密顿量（不必欠扭），
我们把下述逼近结果与 $\rho(\mathord\cdot;1)$ 的 $C^0$-连续性结合使用。厄舍在
(Ush10a) 定理 4.5 中证明了这一结果；它改进了 (Sch06) 定理 1.3，并在逼近过程中
增补了非退化性。

\begin{theorem}
\label{thm:22-slow-approximation-preview}
设 $H:M\to\mathbb R$ 是任意慢哈密顿量。则对任意 $\epsilon>0$，存在光滑函数
$K:M\to\mathbb R$，使 $\lVert K-H\rVert_{C^0}<\epsilon$，且 $K$ 是慢的、欠扭的
莫尔斯函数。
\end{theorem}

\begin{proof}
证明见下一小节，即第 22.2.5 小节。
\end{proof}

最后总结极小性猜想的证明。

\begin{theorem}[厄舍 (Ush10a)]
\label{thm:22-usher-minimality}
设 $(M,\omega)$ 是任意闭辛流形，$G$ 是任意慢自治哈密顿量。则单参数群
$\phi_G^t$ 对 $0\leq t\leq1$ 在其端点固定的同伦类中是长度极小的。
\end{theorem}

\begin{proof}
这是命题~\ref{prop:22-negative-length-minimizing} 与定理
~\ref{thm:22-slow-fundamental-cycle-tight} 的直接推论。更精确地，定理
~\ref{thm:22-slow-fundamental-cycle-tight} 蕴含 $\rho(G;1)=E^-(G)$，再由命题
~\ref{prop:22-negative-length-minimizing} 得 $G$ 使负霍费尔长度极小。因为 $G$
慢时 $\overline G$ 也慢，还有
$\rho(\overline G;1)=E^-(\overline G)=E^+(G)$。因此 $\overline G$ 也使负霍费尔
长度极小，这等价于 $G$ 使正长度极小。两者结合显然说明 $G$ 使霍费尔长度极小。
\end{proof}

\subsection{慢哈密顿量的逼近定理}

本小节旨在逐字依照 (Ush10a) 中厄舍的证明，给出定理
~\ref{thm:22-slow-approximation-preview} 的证明。

\stfhIndex{I200}{Yueke}{约克（Yorke）}
先从约克 (Yo69) 得到的一般常微分方程周期估计开始。设
$F:\Omega\to\mathbb R^n$ 是区域 $\Omega\subset\mathbb R^n$ 上的利普希茨
向量场，以 $\lVert\mathord\cdot\rVert$ 表示欧氏范数。约克 (Yo69)
证明了如下结果。（证明细节请读者参阅原始论文。）

\begin{theorem}
\label{thm:22-yorke-period-estimate}
设 $F$ 是利普希茨常数为 $L>0$ 的利普希茨映射，即
\begin{equation}
 \lVert F(x_1)-F(x_2)\rVert\leq L\lVert x_1-x_2\rVert
 \qquad\text{对每个 }x_1,x_2\in\Omega.
 \label{eq:22-lipschitz-vector-field}
\end{equation}
若 $x=x(t)$ 是 $\dot x=F(x)$ 的周期为 $p$ 的非常值周期解，则
$p\geq2\pi/L$。
\end{theorem}

涩谷（未发表工作）更早证明过一条类似定理，所得较弱结论为 $p\geq2/L$；对这里的
目的而言，这其实已经足够。

这里再次回忆定理~\ref{thm:22-slow-approximation-preview} 的陈述。

\begin{theorem}
\label{thm:22-slow-approximation}
设 $H:M\to\mathbb R$ 是任意慢哈密顿量。则对任意 $\epsilon>0$，存在光滑函数
$K:M\to\mathbb R$，使 $\lVert K-H\rVert_{C^0}<\epsilon$，且 $K$ 是慢的、通有
欠扭的莫尔斯函数。
\footnote{译者注：原书定理 22.2.17 的预告版本称逼近函数为“欠扭”，本定理重述时则称
“通有欠扭”；译文分别照录。}
\end{theorem}

\begin{proof}
选择由与 $\omega$ 相容的殆复结构诱导的 $M$ 上背景黎曼度量，使得
$X_H=J\nabla H$，特别地，对所有 $p$ 都有 $|X_H(p)|=|\nabla H(p)|$。令
\[
 F_n=\{p\in M\mid |\nabla H(p)|\leq1/n\}.
\]
注意，若 $s\in\bigcap_{n=1}^\infty H(F_n)$，则存在 $p_n\in M$ 使
$H(p_n)=s$ 且 $|\nabla H(p_n)|\leq1/n$；由 $M$ 的紧性以及 $H,\nabla H$ 的连续性，
存在 $p\in M$ 满足 $H(p)=s$ 且 $\nabla H(p)=0$。故
$\bigcap_{n=1}^\infty H(F_n)$ 恰为 $H$ 的临界值集合，由萨德定理它的勒贝格
测度为零。因为 $H$ 的像有有限勒贝格测度，故当 $n\to\infty$ 时
$m(H(F_n))\to0$。

给定 $\zeta>0$。存在 $N$ 使 $m(H(F_N))<\zeta/2$。选择光滑函数
$\psi:M\to\mathbb R$，使
\[
 0\leq\psi\leq1,\qquad \psi|_{F_{2N}}=1,
 \qquad\operatorname{supp}(\psi)\subset F_N.
\]
因为 $F_N$ 从而 $H(F_N)\subset\mathbb R$ 紧，所以存在闭包两两不交的开区间
$I_1,\ldots,I_k$，使
\[
 H(F_N)\subset\bigcup_{j=1}^k I_j,
 \qquad\sum_{j=1}^k m(I_j)<\frac{3\zeta}{4}.
\]
再选闭包两两不交的开区间 $I'_1,\ldots,I'_k$，使对每个 $j$ 都有
$I_j\Subset I'_j$ 且 $\sum_{j=1}^k m(I'_j)<\zeta$。记
\[
 S=\bigcup_{j=1}^k I_j,
 \qquad S'=\bigcup_{j=1}^k I'_j.
\]
令
\[
 B=\sup_{p\in M}|\nabla(\nabla H)(p)|,
\]
并取一个稍后还要进一步指定的小数 $\eta>0$；特别要求 $\eta<1$。选择
$C^\infty$ 函数
$f:[\min H,\max H]\to[\min H,\max H]$ 满足：
\begin{enumerate}
 \item $f(\min H)=\min H$；
 \item 对所有 $s$，$0<f'(s)\leq1$；
 \item 若 $s\notin S'$，则 $f'(s)=1$；
 \item 若 $s\in S$，则 $f'(s)=\eta/2$。
\end{enumerate}
因为 $m(S')<\zeta$，对所有 $s\in[\min H,\max H]$ 都有
$0\leq s-f(s)<\zeta$。特别地，$\lVert f\circ H-H\rVert_{C^0}<\zeta$，且此估计
与 $\eta$ 无关。

若 $q\in F_N$（由 $f$ 的构造，这特别蕴含 $f''(H(q))=0$），则
\[
 |\nabla(f\circ H)(q)|=f'(H(q))|\nabla H(q)|\leq\frac{\eta}{2N}
\]
且
\[
 |\nabla(\nabla(f\circ H))(q)|
 =\frac{\eta}{2}|\nabla(\nabla H)(q)|\leq\frac{B\eta}{2}.
\]
故
\[
 \lVert X_{f\circ H}\rVert_{C^1}
 \leq\frac{\eta}{2}+\frac{B\eta}{2}.
\]
回忆上述函数 $\psi:M\to\mathbb R$（其支集包含于 $F_N$），于是
\footnote{译者注：此处 $\psi X_{f\circ H}$ 表示函数 $\psi$ 与向量场
$X_{f\circ H}$ 的逐点乘积；源书即印作并列乘积，无复合圈。}
\[
 \lVert\psi X_{f\circ H}\rVert_{C^1}
 \leq\left(\frac{\eta}{2}+\frac{B\eta}{2}\right)
 \lVert\psi\rVert_{C^1}.
\]
由定理~\ref{thm:22-yorke-period-estimate} 的估计（以及惠特尼嵌入定理），存在
$\eta_0>0$，使 $M$ 上任意满足
$\lVert V\rVert_{C^1}\leq\eta_0(B+1)\lVert\psi\rVert_{C^1}$ 的向量场 $V$
都没有周期小于或等于一的非常值周期轨道。

于是取 $0<\eta<\eta_0$，再取 $M$ 上一列莫尔斯函数
$\{K_m\}_{m=1}^\infty$，使 $K_m$ 在 $C^2$-范数下趋于 $f\circ H$。则
$X_{K_m}\to X_{f\circ H}$ 于 $C^1$-范数，因而
$\psi X_{K_m}\to\psi X_{f\circ H}$ 于 $C^1$-范数。所以当 $m$ 足够大，使
\[
 \lVert\psi X_{f\circ H}-\psi X_{K_m}\rVert_{C^1}
 <\frac{(B+1)\eta}{2}\lVert\psi\rVert_{C^1},
\]
向量场 $\psi X_{K_m}$ 没有周期小于或等于一的非常值周期轨道。

现在，因为 $H$ 慢且 $0<f'\leq1$，容易看出 $f\circ H$ 也慢。因此
$X_{f\circ H}$ 的所有周期小于或等于一的可缩周期轨道，都是 $f\circ H$ 临界点处
的常值轨道；特别地，它们全都包含在 $F_{2N}$ 的内部。阿尔泽拉--阿斯科利定理的一个
简单应用于是表明：若 $m$ 充分大，$X_{K_m}$ 的任意周期小于或等于一的可缩周期轨道
必包含在 $F_{2N}$ 中。因此（由于 $\psi|_{F_{2N}}=1$），对充分大的 $m$，
$X_{K_m}$ 的周期小于或等于一的可缩周期轨道其实也是 $\psi X_{K_m}$ 的周期轨道，
从而由上文可知只能是常值轨道。所以充分大的 $m$ 使 $K_m$ 为慢的。

此外，$K_m\to f\circ H$ 于 $C^2$ 拓扑蕴含：当 $m$ 充分大时，$K_m$ 的所有临界点
$p$ 都包含于 $F_N$；在任一这样的 $p$ 处，
\[
 |\nabla(\nabla(f\circ H))(p)|<B\eta/2,
\]
所以当 $m$ 充分大时 $|\nabla(\nabla K_m)(p)|<B\eta$。因此，只要起初把 $\eta$
取成小于 $2\pi/B$，$K_m$ 在任一此类临界点处的线性化流 $\phi_{K_m}^t$ 就没有周期
小于或等于 $1$ 的非常值周期轨道。于是对充分大的 $m_0$，$K_{m_0}$ 是慢的、平坦的
莫尔斯函数，且
$\lVert K_{m_0}-f\circ H\rVert_{C^2}<\eta$。令 $\zeta=\epsilon/2$，并取
$K=K_{m_0}$，再结合 $\lVert f\circ H-H\rVert_{C^0}<\eta$，即证定理。
\footnote{译者注：原证明取 $\zeta=\epsilon/2$，但对
$\lVert K_{m_0}-f\circ H\rVert_{C^2}<\eta$ 及
$\lVert f\circ H-H\rVert_{C^0}<\eta$ 未再明确要求 $\eta<\epsilon/2$；此外，此处结论
称 $K_{m_0}$ “平坦”，而定理陈述称“通有欠扭”。
译文保留论证结构并提示这两处缺口。}
\end{proof}

\section{谱容量与最优能量--容量不等式}
\label{sec:spectral-capacities-energy-capacity}

本节研究若干已知辛容量构造的各种谱类比。我们以霍费尔直径、霍费尔位移能以及
$\pi_1$-敏感霍费尔--曾德尔容量为例。谱位移能由 (Oh05d) 引入；
$\pi_1$-敏感霍费尔--曾德尔容量由\personLu{} (LuG96) 与施瓦茨 (Schw00) 引入，
随后厄舍 (Ush10a) 用谱不变量进一步研究了它。

\stfhIndex{I173}{Puzhijing}{谱直径（spectral diameter）}
\begin{definition}[谱直径]
\label{def:22-spectral-diameter}
设 $(M,\omega)$ 是不必无边界的紧辛流形。考虑支集包含于 $M\setminus\partial M$ 的
哈密顿微分同胚群 $\operatorname{Ham}_c(M,\omega)$。定义
$\operatorname{Ham}_c(M,\omega)$ 的\emph{谱直径}为
\[
 \begin{aligned}
 \operatorname{diam}_\rho(\operatorname{Ham}_c(M,\omega))
 &:=\sup_H\{\gamma(H)\mid\operatorname{supp}H\subset M\setminus\partial M\}\\
 &=\sup_{\phi\in\operatorname{Ham}_c(M,\omega)}\{\gamma(\phi)\}.
 \end{aligned}
\]
\end{definition}

一般而言，尚不清楚哪些辛流形 $(M,\omega)$ 满足
$\operatorname{diam}_\rho(M,\omega)<\infty$。稍后将在第 22.5 节证明
$\mathbb{CP}^n$ 的谱直径有限。

然而，若考虑可位移开子集 $U\subset M$，便可用\emph{谱位移能}（从而也可用霍费尔位移
能量）给出 $\operatorname{diam}_\rho(\operatorname{Ham}_c(U,\omega))$ 的一个界。
\stfhIndex{I174}{Puweiyineng}{谱位移能（spectral displacement energy）}
现在引入谱位移能。

\begin{definition}
\label{def:22-spectral-displacement-energy}
设 $A$ 是 $M$ 的任意闭子集。定义 $A$ 的谱位移能 $e_\rho(A)$ 为
\begin{equation}
 e_\rho(A)=\inf_H\{\gamma(H)\mid\phi_H^1(A)\cap A=\varnothing\}.
 \label{eq:22-spectral-displacement-energy}
\end{equation}
\end{definition}

为简化记号，以
\[
 \mathcal P_U^{\mathrm{ham}}
 :=\mathcal P^{\mathrm{ham}}(\operatorname{Symp}_U(M,\omega),\mathrm{id})
\]
表示所有满足 $\operatorname{supp}F\subset U$ 的哈密顿路径 $\phi_F$ 之集。

有了这些定义，我们陈述如下结论。

\begin{proposition}[与 (EnP06) 引理 7.2 比较]
\label{prop:22-displaceable-spectral-diameter}
设 $U$ 是闭包 $\overline U$ 可位移的开子集。对任意
$\lambda=\phi_F\in\mathcal P_U^{\mathrm{ham}}$，都有
\[
 \gamma(\phi_F)\leq2e_\rho(\overline U).
\]
特别地，
\[
 \operatorname{diam}_\rho(\operatorname{Ham}_c(U))\leq2e_\rho(\overline U).
\]
\end{proposition}

\begin{proof}
设 $H$ 是满足 $\phi_H^1(U)\cap\overline U=\varnothing$ 的任意哈密顿量，$F$ 是满足
$\operatorname{supp}F\subset U$ 的哈密顿量。
\footnote{译者注：原书此处连续两次把新引入的哈密顿量写作 $F$，但首个对象满足的位移
条件写作 $\phi_H^1(U)\cap\overline U=\varnothing$；译文依该条件把首个对象改作 $H$。}
注意，$\gamma(\lambda)$ 的值不依赖于
代表 $\lambda=\phi_F$ 的哈密顿量 $F$ 的归一化。因此，为证明命题，只需对满足
$\operatorname{supp}F\subset U$、但不必均值归一化的哈密顿量 $F$ 证明
\begin{equation}
 \gamma(F)=\rho(\phi_F;1)+\rho(\phi_F^{-1};1)\leq2\gamma(H).
 \label{eq:22-local-gamma-bound}
\end{equation}
先从如下引理开始。

\begin{lemma}
\label{lem:22-calabi-spectral-shift}
有
\[
 \begin{aligned}
 \rho(\phi_H\phi_F;1)&=\rho(\phi_H;1)+\operatorname{Cal}_U(F),\\
 \rho(\phi_H^{-1}(\phi_F)^{-1};1)
 &=\rho(\phi_H^{-1};1)-\operatorname{Cal}_U(F).
 \end{aligned}
\]
\end{lemma}

\begin{proof}
证明第一式即可，因为 $(\phi_F)^{-1}=\phi_{\overline F}$、
$\operatorname{supp}\overline F\subset U$，并且当 $\operatorname{supp}F\subset U$ 时
$\operatorname{Cal}^{\mathrm{path}}(\phi_{\overline F})
=-\operatorname{Cal}^{\mathrm{path}}(\phi_F)$。我们依照 (EnP06) 引理 7.2 的方案；该方案
\stfhIndex{I148}{Aositeluofu}{奥斯特罗弗（Ostrover）}
又基于奥斯特罗弗 (Os03) 的一个技巧。

对 $0\leq s\leq1$，考虑族 $\lambda_s$：
\[
 \lambda_s(t)=
 \begin{cases}
  \phi_H^{2t},&0\leq t\leq\frac12,\\
  \phi_H^1\phi_F^{2s(t-1/2)},&\frac12\leq t\leq1.
 \end{cases}
\]
有
\[
 \operatorname{Dev}(\lambda_s)(t)=
 \begin{cases}
  2H(2t,x),&0\leq t\leq\frac12,\\
  2sF(2s(t-1/2),x),&\frac12\leq t\leq1.
 \end{cases}
\]
记 $F^s:=\operatorname{Dev}(\lambda_s)$。注意
$\operatorname{Fix}(\phi_H^1\lambda_s)=\operatorname{Fix}(\phi_H^1)$。对每个不动点
$x\in\operatorname{Fix}(\phi_H^1\lambda_s)\equiv\operatorname{Fix}\phi_H^1$，以
$z_x^s$ 表示 $F^s$ 的相应周期轨道。则当 $0\leq t\leq1/2$ 时
$z_x^s(t)=\phi_H^{2t}(x)$，当 $1/2\leq t\leq1$ 时
$z_x^s(t)\equiv\phi_H^1(x)=x$。因此可取光滑提升
$[z_x^s,w]\in\widetilde\Omega_0(M)$，其中填充圆盘 $w$ 与 $s$ 无关。于是对每个
$s\in[0,1]$，
\[
 \begin{aligned}
 \mathcal A_{F^s}([z_x^s,w])
 &=-\int w^*\omega-\int_0^1F^s(t,z_x^s(t))\,dt\\
 &=-\int w^*\omega
 -\int_0^{1/2}2H(2t,\phi_H^{2t}(x))\,dt\\
 &\quad-\int_{1/2}^1 2sF(2s(t-1/2),\phi_H^1(x))\,dt.
 \end{aligned}
\]
注意
\[
 \int_{1/2}^1F^s(t,z_x^s(t))\,dt
 =\int_{1/2}^1 2sF(2s(t-1/2),\phi_H^1(x))\,dt
 =\int_0^sF(t,x)\,dt,
\]
其中第二个等号来自 $x\in\operatorname{Fix}\phi_H^1$；并且
\[
 \int_0^{1/2}2H(2t,\phi_H^{2t}(x))\,dt
 =\int_0^1H(t,\phi_H^t(x))\,dt.
\]
由此得到
\[
 \mathcal A_{F^s}([z_x^s,w])
 =\mathcal A_H([z_x^0,w])-\int_0^sF(t,x)\,dt.
\]
但对任意 $x\in M\setminus U$，
\begin{equation}
 F(t,x)\equiv\{\operatorname{vol}(M)\}^{-1}\int_MF_t\,d\mu.
 \label{eq:22-hamiltonian-outside-support}
\end{equation}
所以对所有 $x\in\operatorname{Fix}\phi_H^1$ 与 $s\in[0,1]$，
\[
 \mathcal A_{F^s}([z_x^s,w])
 =\mathcal A_H([z_x^0,w])
 -\int_0^s\{\operatorname{vol}(M)\}^{-1}\int_MF_t\,d\mu\,dt,
\]
从而
\[
 \mathcal A_{F^s}([z_x^s,w])\in\operatorname{Spec}H+c(s),
 \qquad
 c(s)=\int_0^s\operatorname{vol}_\omega(M)^{-1}\int_MF_t\,d\mu\,dt.
\]
特别地，作为 $\mathbb R$ 的子集，有
\footnote{译者注：原书上一行把作用量写成属于
$\operatorname{Spec}H+c(s)$，此式却给出
$\operatorname{Spec}F^s=\operatorname{Spec}H-c(s)$；译文逐式照录并提示符号不一致。}
\begin{equation}
 \operatorname{Spec}F^s=\operatorname{Spec}H-c(s).
 \label{eq:22-spectrum-shift}
\end{equation}

现考虑连续函数 $s\mapsto\rho(\lambda_s;1)=\rho(F^s;1)$。由于所有 $F^s$ 都非退化，
\footnote{译者注：命题 22.3.3 并未假设 $H$ 非退化，此处也没有给出逼近步骤，却直接
断言所有 $F^s$ 非退化；译文照录。}
非退化谱性公理给出 $\rho(F^s;1)\in\operatorname{Spec}F^s$，故由
~\eqref{eq:22-spectrum-shift}，
\[
 \rho(F^s;1)+c(s)\in\operatorname{Spec}H.
\]
因为 $s\mapsto\rho(F^s;1)+c(s)$ 连续，而 $\operatorname{Spec}H$ 无处稠密，前者必为
常值。特别地，
\[
 \rho(F^0;1)+c(0)=\rho(F^1;1)+c(1).
\]
但 $c(0)=0$ 且
\[
 c(1)=\frac1{\operatorname{vol}_\omega(M)}
 \int_0^1\int_MF\,d\mu\,dt=\operatorname{Cal}_U(F),
\]
又有 $F^0=H$、$F^1=H*F$。所以
\[
 \rho(\phi_H;1)=\rho(\phi_H\phi_F;1)-\operatorname{Cal}_U(F),
\]
等价地，
\begin{equation}
 \rho(\phi_H\phi_F;1)=\rho(\phi_H;1)+\operatorname{Cal}_U(F).
 \label{eq:22-calabi-shift}
\end{equation}
引理得证。
\end{proof}


回到~\eqref{eq:22-local-gamma-bound} 与命题
~\ref{prop:22-displaceable-spectral-diameter} 的证明。把引理
~\ref{lem:22-calabi-spectral-shift} 中两式相加，得到
\begin{equation}
 \rho(\phi_H\phi_F;1)+\rho(\phi_H^{-1}\phi_F^{-1};1)
 =\rho(\phi_H;1)+\rho(\phi_H^{-1};1).
 \label{eq:22-calabi-sum}
\end{equation}
由三角不等式，还有
\[
 \begin{aligned}
 \rho(\phi_F;1)&\leq\rho(\phi_H\phi_F;1)+\rho(\phi_H^{-1};1),\\
 \rho(\phi_F^{-1};1)&\leq
 \rho(\phi_F^{-1}\phi_H^{-1};1)+\rho(\phi_H;1).
 \end{aligned}
\]
把这两个不等式相加，并将所得和代入~\eqref{eq:22-calabi-sum}，便得到
\[
 \gamma(F)\leq2\rho(\phi_H;1)+2\rho(\phi_H^{-1};1)=2\gamma(H),
\]
从而完成证明。
\end{proof}

此时尚不显然 $\operatorname{diam}_\rho(\operatorname{Ham}_c(U))>0$。依照
(LuG96), (Schw00), (Ush10a) 的论证，我们将通过把 $\rho(\mathord\cdot;1)$ 与另一
个称为 $\pi_1$-敏感霍费尔--曾德尔容量的容量比较来证明这一点。这一概念最初由\personLu{} (LuG96)
引入；这里采用施瓦茨在 (Schw00) 中使用的术语。

为引出卢广存对 $\pi_1$-敏感霍费尔--曾德尔容量的定义，先回忆霍费尔--曾德尔
容量 $c_{\mathrm{HZ}}(A)$ 的构造。为定义 $c_{\mathrm{HZ}}(A)$，回忆\emph{假设 $M$ 有非空
边界} $\partial M$。记
\[
 \begin{aligned}
 \mathcal H_{\mathrm{HZ}}(A)=\{H\in C^\infty(M)\mid{}&
 \operatorname{supp}(H)\subset A\setminus\partial M,
 \quad0\leq H\leq\max H,\\
 &H^{-1}(0)\text{ 与 }H^{-1}(\max H)
 \text{ 均包含非空开集}\}.
 \end{aligned}
\]

\stfhIndex{I004}{00-hz-rongxu}{HZ-容许}
\begin{definition}
\label{def:22-hz-admissible}
若 $H\in\mathcal H_{\mathrm{HZ}}(A)$ 作为一周期函数时，其哈密顿流 $\phi_H$ 没有周期
小于或等于 $1$ 的非常值周期轨道，则称 $H$ 为 \emph{HZ-容许的}。类似地，若
$\phi_H$ 没有在 $M$ 中可缩、周期小于或等于 $1$ 的非常值周期轨道，则称 $H$ 为
\emph{$\mathrm{HZ}^\circ$-容许的}。分别以
$\mathcal H_{\mathrm{HZ}}^{\mathrm{ad}}(A)$ 和
$\mathcal H_{\mathrm{HZ}}^{\circ,\mathrm{ad}}(A,M)$ 表示 HZ-容许与
$\mathrm{HZ}^\circ$-容许哈密顿量的集合。（路使用 C-容许性一词，而非
$\mathrm{HZ}^\circ$-容许性。）
\end{definition}

霍费尔--曾德尔容量 (HZ94) 定义如下。

\stfhIndex{I113}{Huofeier-Zengdeerrongliang}{霍费尔--曾德尔容量（Hofer--Zehnder capacity）}
\begin{definition}[霍费尔--曾德尔容量]
\label{def:22-hofer-zehnder-capacity}
设 $A\subset M\setminus\partial M$。定义
\footnote{译者注：原书此处印作 $A\subset M\subset\partial M$，与“$M$ 有非空边界”
以及上文 $\operatorname{supp}(H)\subset A\setminus\partial M$ 不相容；译文据定义域改作
$A\subset M\setminus\partial M$。}
\[
 c_{\mathrm{HZ}}(A)=\sup\{\max H\mid
 H\in\mathcal H_{\mathrm{HZ}}(A)\text{ 是 HZ-容许的}\},
\]
称之为霍费尔--曾德尔容量。
\end{definition}

\stfhIndex{I172}{Shiwaci}{施瓦茨（Schwarz）}
把 $\pi_1(M)$ 对哈密顿流的影响纳入考虑后，路 (LuG96) 与施瓦茨 (Schw00) 引入
$c_{\mathrm{HZ}}$ 的如下变体：
\[
 c_{\mathrm{HZ}}^\circ(A,M)
 :=\sup\{\max H\mid H\in\mathcal H_{\mathrm{HZ}}(A)
 \text{ 是 $\mathrm{HZ}^\circ$-容许的}\}.
\]
他们称之为 $\pi_1$-敏感霍费尔--曾德尔容量。在该定义中，$M$ \emph{现在可以是闭的}。

这两个记号符合如下事实：$c_{\mathrm{HZ}}(A)$ 只依赖于
$(A,\phi_H^t|_A)$（以及 $\dim M$）；而 $\phi_H^t$ 的轨道在 $M$ 中是否可缩依赖于
$M$，故 $c_{\mathrm{HZ}}^\circ(A,M)$ 依赖于环境流形 $M$，而不仅依赖于 $A$。
显然，任意 HZ-容许哈密顿量都是 $\mathrm{HZ}^\circ$-容许的，所以
\[
 c_{\mathrm{HZ}}(A)\leq c_{\mathrm{HZ}}^\circ(A,M).
\]
霍费尔--曾德尔在 (HZ94) 中证明 $c_{\mathrm{HZ}}$ 是辛容量；若
$B^{2n}(r)$ 是 $\mathbb R^{2n}$ 中半径为 $r$ 的球，则
$c_{\mathrm{HZ}}(B^{2n}(r))=\pi r^2$。

\stfhIndex{I190}{Eshe}{厄舍（Usher）}
下述结果由厄舍 (Ush10a) 证明。

\begin{proposition}
\label{prop:22-displaced-support-spectral-bound}
设 $F,H:(\mathbb R/\mathbb Z)\times M\to\mathbb R$ 是无边界紧辛流形
$(M,\omega)$ 上两个不必归一化的哈密顿量，并满足
\[
 \phi_H^1(\operatorname{supp}F)\cap\operatorname{supp}F=\varnothing,
\]
其中
\[
 \operatorname{supp}F:=\bigcup_{t\in\mathbb R/\mathbb Z}
 \operatorname{supp}F_t\subset M.
\]
则
\[
 \rho(F;1)\leq\gamma(H)
 =\rho(H;1)+\rho(\overline H;1).
\]
\end{proposition}

\begin{proof}
选择可位移开子集 $U\subset M$，使 $\operatorname{supp}F\subset U$。记
$\underline F$ 为 $F$ 的归一化：
\[
 \underline F_t=F_t-\frac1{\operatorname{vol}(M)}\int_MF_t\,\omega^n
\]
对每个 $t$ 成立。由定义，
\begin{align}
 \rho(F;1)
 &=\rho(\underline F;1)+\frac1{\operatorname{vol}(M)}
 \int_0^1\int_MF_t\,\omega^n\,dt\notag\\
 &=\rho(\underline F;1)+\operatorname{Cal}_U(F)\notag\\
 &\leq\rho(F\#H;1)+\rho(\overline H;1)+\operatorname{Cal}_U(F).
 \label{eq:22-displaced-support-intermediate}
\end{align}
\footnote{译者注：原书式 (22.3.30) 第一行的空间积分 $\int_MF_t$ 未写体积形式；同页
归一化公式明确使用 $\omega^n$，译文据此补出。}
另一方面，~\eqref{eq:22-calabi-shift} 蕴含：对每个满足
$\phi_H^1(U)\cap\overline U=\varnothing$ 的 $H$，
\[
 \operatorname{Cal}_U^{\mathrm{path}}(\phi_F)
 =\rho(H;1)-\rho(F\#H;1).
\]
把它代入~\eqref{eq:22-displaced-support-intermediate}，便得到
\[
 \rho(F;1)\leq\rho(H;1)+\rho(\overline H;1).
\]
\end{proof}

结合定理~\ref{thm:22-slow-fundamental-cycle-tight} 与此命题，得到如下定理。

\begin{theorem}[厄舍]
\label{thm:22-pi-one-capacity-displacement}
若 $(M,\omega)$ 闭，则对任意 $A\subset M$ 有
\[
 c_{\mathrm{HZ}}^\circ(A,M)\leq e(A,M).
\]
\end{theorem}

\begin{proof}
以一列慢、欠扭且非退化的自治哈密顿量 $F_i$ 逼近任意容许的 $F$。由定理
~\ref{thm:22-slow-fundamental-cycle-tight}，
$\rho(-F_i;1)=E^+(F_i)$。利用 $\rho(\mathord\cdot;1)$ 与 $E^+$ 的 $C^0$-连续性取
极限，得到 $\rho(-F;1)=E^+(F)=\max F$。于是命题
~\ref{prop:22-displaced-support-spectral-bound} 蕴含
\[
 \max F\leq\rho(H;1)+\rho(\overline H;1).
\]
让右端对所有满足 $\phi_H^1(U)\cap\overline U=\varnothing$ 的 $H$ 取下确界，再让
$\max F$ 对所有 $\mathrm{HZ}^\circ$-容许的 $F$ 取上确界，即得定理。
\end{proof}

定理~\ref{thm:22-pi-one-capacity-displacement} 的一个直接推论是谱范数的非退化性。

\begin{corollary}
\label{cor:22-spectral-norm-nondegenerate}
谱范数 $\gamma:\operatorname{Ham}(M,\omega)\to\mathbb R_+$ 满足
$\gamma(\phi)=0$ 当且仅当 $\phi=\mathrm{id}$。
\end{corollary}

\begin{proof}
设 $\phi\neq\mathrm{id}$。则可找到一个可位移的充分小辛球 $B\cong B^{2n}(r)$。
\footnote{译者注：原书这里只说所选小球“可位移”，但随即用它给出给定 $\phi$ 的谱范数
下界；论证还需选择满足 $\phi(B)\cap B=\varnothing$ 的充分小球。}
于是 $\gamma(\phi)\geq c_{\mathrm{HZ}}^\circ(B,M)$。另一方面，
\footnote{译者注：定义 22.3.6 中 $c_{\mathrm{HZ}}$ 只有集合参数，此处原书却写
$c_{\mathrm{HZ}}(B,M)$；译文保留二参数写法。}
\[
 c_{\mathrm{HZ}}^\circ(B,M)\geq c_{\mathrm{HZ}}(B,M)=\pi r^2>0.
\]
结合这两个不等式便有 $\gamma(\phi)>0$。
\end{proof}

另一推论给出格罗莫夫面积与霍费尔位移能之间如下的最优不等式，即能量--容量
不等式；证明见 (Ush10a), (Sch06)。

\begin{corollary}[最优能量--容量不等式]
\label{cor:22-sharp-energy-capacity}
设 $A\subset M$ 是任意紧闭子集，以 $c(A)$ 表示格罗莫夫面积，以 $e(A)$ 表示霍费尔
位移能。则
\footnote{译者注：原书在本推论中先以 $e(A)$ 表示霍费尔位移能，结论却改写为
$e(A,M)$；译文逐处照录这一参数数目不一致。}
\[
 c(A)\leq e(A,M).
\]
\end{corollary}

\begin{proof}
众所周知（见命题 6.1.5），由于格罗莫夫面积是最小的可能辛容量，
$c(A)\leq c_{\mathrm{HZ}}(A)$。因此，定理
~\ref{thm:22-pi-one-capacity-displacement} 以及不等式
$c_{\mathrm{HZ}}(A)\leq c_{\mathrm{HZ}}^\circ(A,M)$ 立即给出此推论。
\end{proof}

\section{恩托夫和波尔捷罗维奇的部分辛拟态}

本节从谱不变量 $\rho(\,\cdot\,;1)$（其中 $1=PD[M]$）中提取若干辛拓扑性质。
\stfhIndex{I074}{Entuofu-Boerjieluoweiqi}{恩托夫--波尔捷罗维奇（Entov--Polterovich）}
为此，我们紧随恩托夫和波尔捷罗维奇 (EnP06) 所采用的步骤。

考虑渐近平均
\begin{equation}
 \mu_1^{\mathrm{path}}(\phi_H)
 =\lim_{k\to\infty}\frac{\rho(\phi_H^k;1)}{k}.
 \label{eq:22-path-asymptotic-average}
\end{equation}
\emph{一旦}证明该定义良定，依定义便有
\begin{equation}
 \mu_1^{\mathrm{path}}(\phi_H^m)=m\mu_1^{\mathrm{path}}(\phi_H).
 \label{eq:22-path-homogeneity}
\end{equation}
我们将在下文命题~\ref{prop:22-path-quasimorphism} 中证明其良定性。

令 $F\in C^\infty(M)$ 为任意自治函数，不要求归一化。考虑它的归一化
\[
 \underline F=F-\frac{1}{\operatorname{vol}_\omega(M)}\int_M F\,d\mu_\omega.
\]
于是恩托夫和波尔捷罗维奇考虑极限
\begin{equation}
 \zeta_{\mathrm{sp}}(F):=\lim_{k\to\infty}\frac{\rho(k\underline F;1)}{k},
 \label{eq:22-spectral-quasi-state}
\end{equation}
其中下标 ``sp'' 代表 ``spectral''。回忆定义 21.3.9 中
$\rho(\phi_{\underline H};1):=\rho(\underline H;1)$ 的定义。由于
\[
 F=\underline F+\frac{1}{\operatorname{vol}(M)}\int_M F\,d\mu_\omega,
\]
定义~\eqref{eq:22-spectral-quasi-state} 可改写为
\begin{align*}
 \zeta_{\mathrm{sp}}(F)
 &=\lim_{k\to\infty}\frac{\rho(k\underline F;1)}{k}
   +\frac{\int_M F\,d\mu_\omega}{\operatorname{vol}(M)}\\
 &=\lim_{k\to\infty}\frac{\rho(\phi_F^k;1)}{k}
   +\frac{\int_M F\,d\mu_\omega}{\operatorname{vol}(M)},
\end{align*}
第二个等号成立是因为 $F$ 自治。\footnote{译者注：按
\eqref{eq:22-spectral-quasi-state} 的定义
$\zeta_{\mathrm{sp}}(F)=\lim_{k\to\infty}\rho(k\underline F;1)/k$，上式第一个
等号右端不应含加项 $\frac{\int_MF\,d\mu_\omega}{\operatorname{vol}(M)}$，即第一个
等号按定义不成立；源书本意应为
$\zeta_{\mathrm{sp}}(F)=\lim_{k\to\infty}\frac{\rho(k\underline F;1)}{k}
=\lim_{k\to\infty}\frac{\rho(\phi_F^k;1)}{k}
+\frac{\int_MF\,d\mu_\omega}{\operatorname{vol}(M)}$。此外，本节末尾命题
22.4.10 证明的收尾（以卡拉比性质代入）只在平移规则 $\rho(H+c)=\rho(H)+c$ 下
闭合，而在该约定下本式的两个加号应改为负号；反之按 $\rho(H+c)=\rho(H)-c$ 则
收尾残留 $+2\frac{\int_MH\,d\mu_\omega}{\operatorname{vol}(M)}$ 而不能闭合。
即源书的重排式符号、卡拉比值符号与收尾结论在任一自洽平移约定下无法同时成立。
结论 $\zeta_{\mathrm{sp}}(F+H)=\zeta_{\mathrm{sp}}(F)$ 本身是标准真命题（等价于
$\zeta_{\mathrm{sp}}$ 在支集可移的函数上恒为零）。译文照录原书符号，特此披露。}

可以把 $\zeta_{\mathrm{sp}}$ 看作 $\mu_1^{\mathrm{path}}$ 的“无穷小”版本。回忆
$\Ham(M,\omega)$ 在恒等元处的曲线芽对应其在恒等元处的切空间，后者同构于归一化
自治哈密顿量的集合。由 $\rho$ 的 $C^0$-连续性，$\zeta_{\mathrm{sp}}$ 连续延拓到
$C^0(M)$。

恩托夫和波尔捷罗维奇在 (EnP06) 中观察到，$\zeta_{\mathrm{sp}}$ 的行为类似阿尔内斯
(Aa91) 引入的\emph{拟态}，并提出了上述 $\zeta_{\mathrm{sp}}$ 的定义。按 (EnP06) 的术语，
$\zeta_{\mathrm{sp}}$ 是\emph{部分辛拟态}的一个例子。

先从阿尔内斯的一般拟态概念开始。

\stfhIndex{I025}{Aerneisi}{阿尔内斯（Aarnes）}
\begin{definition}[阿尔内斯]
\label{def:22-aarnes-quasi-state}
令 $M$ 为拓扑空间，以 $C^0(M)$ 表示 $M$ 上连续实值函数的集合。若（非线性）泛函
\stfhIndex{I161}{Nitai}{拟态（quasi-state）}
$\zeta:C^0(M)\to\mathbb R$ 满足下列性质，则称它为一个拟态：
\begin{enumerate}[label=\arabic*.]
\item （线性）$\zeta$ 在 $\mathcal A_F$ 上线性，其中 $\mathcal A_F$ 是由
      $F\in C(M)$ 生成的 $\mathbb R$-代数；
\item （单调性）若逐点有 $F\leq G$，则 $\zeta(F)\leq\zeta(G)$；
\item （归一化）$\zeta(1)=1$。
\end{enumerate}
\end{definition}

下面是 (EnP06) 引入的辛版本。

\begin{definition}[恩托夫和波尔捷罗维奇]
\label{def:22-symplectic-quasi-state}
令 $(M,\omega)$ 为辛流形。若（非线性）泛函
$\zeta:C^0(M)\to\mathbb R$ 是定义~\ref{def:22-aarnes-quasi-state}
\stfhIndex{I187}{Xinnitai}{辛拟态（symplectic quasi-state）}
意义下 $M$ 上的拟态，并满足下列附加辛性质，则称它为\emph{辛拟态}：
\begin{enumerate}[label=\arabic*.,start=4]
\item （强拟可加性）设 $F_1,F_2$ 光滑且 $\{F_1,F_2\}=0$，其中
      $\{F_1,F_2\}$ 是与 $\omega$ 相伴的泊松括号，则
      $\zeta(F_1+F_2)=\zeta(F_1)+\zeta(F_2)$；
\item （消失性）若 $\operatorname{supp}F$ 可移，则 $\zeta(F)=0$；
\item （辛不变性）对任意 $\psi\in\Symp_0(M,\omega)$，
      $\zeta(F)=\zeta(F\circ\psi)$，其中 $\Symp_0(M,\omega)$ 是
      $\Symp(M,\omega)$ 的恒等分支。
\end{enumerate}
\end{definition}

事实表明，对一般的 $(M,\omega)$，上述 $\zeta_{\mathrm{sp}}$ 并不完全满足此定义中
（辛）拟态的全部性质，但满足比（辛）拟态较弱的性质。

\begin{definition}[恩托夫--波尔捷罗维奇]
\label{def:22-partial-symplectic-quasi-state}
\emph{部分辛拟态}是指具有下列性质的函数
\[
 \zeta:C^0(M)\longrightarrow\mathbb R.
\]
\begin{enumerate}[label=\arabic*.]
\item （利普希茨连续性）
      $|\zeta(F)-\zeta(H)|\leq\lVert F-H\rVert_{C^0}$；
\item （半齐次性）对任意 $F\in C^0(M)$ 与 $\lambda\in\mathbb R_{\geq0}$，
      $\zeta(\lambda F)=\lambda\zeta(F)$；
\item （单调性）若 $F\leq G$，则 $\zeta(F)\leq\zeta(G)$；
\item （归一化）$\zeta(1)=1$；
\item （部分可加性）若 $F_1,F_2\in C^\infty(M)$ 满足
      $\{F_1,F_2\}=0$ 且 $\operatorname{supp}F_2$ 可移，则
      $\zeta(F_1+F_2)=\zeta(F_1)$；
\item （辛不变性）对所有 $\psi\in\Symp_0(M,\omega)$，
      $\zeta(F)=\zeta(F\circ\psi)$。
\end{enumerate}
\end{definition}

需要指出，定义~\ref{def:22-symplectic-quasi-state} 中的消失性公理是部分可加性公理的
直接推论。有了这个定义，可陈述如下定理。

\begin{theorem}[恩托夫和波尔捷罗维奇 (EnP03)]
\label{thm:22-spectral-partial-quasi-state}
令 $(M,\omega)$ 为任意紧辛流形。函数
$\zeta_{\mathrm{sp}}:C^\infty(M)\to\mathbb R$ 可延拓为函数（仍记作）
\stfhIndex{I150}{Bufenxinnitai}{部分辛拟态（partial symplectic quasi-states）}
$\zeta_{\mathrm{sp}}:C^0(M)\to\mathbb R$，且后者为部分辛拟态。称
\stfhIndex{I177}{Pubufennitai}{谱部分拟态（spectral partial quasi-state）}
$\zeta_{\mathrm{sp}}$ 为谱部分拟态。
\end{theorem}

恩托夫和波尔捷罗维奇陈述该定理时假设 $(M,\omega)$ 强半正且\emph{有理}。不过，一旦证明
非退化谱性公理成立，这些限制便无必要；(Oh09a) 与 (Ush08) 基于用仓西结构
构造虚基本类，在完全一般的情形下证明了这一点。本节末将按照恩托夫和波尔捷罗维奇
借自 (EnP06) 的论证给出证明。

先说明仅凭部分辛拟态的存在便能推出的一些引人注目的辛交性质；这些结果由恩托夫和
波尔捷罗维奇 (EnP06) 证明。

\begin{definition}[恩托夫--波尔捷罗维奇]
\label{def:22-poisson-commutative}
若线性子空间 $\mathcal A\subset C^\infty(M)$ 对所有 $F,G\in\mathcal A$ 都满足
$\{F,G\}=0$，则称 $\mathcal A$ 为泊松可交换的（或对合的）。给定有限维
泊松可交换子空间 $\mathcal A\subset C^\infty(M)$，由方程
\[
 \langle\Phi_{\mathcal A}(x),F\rangle=F(x)
\]
\stfhIndex{I095}{Guangyidongliangyingshe}{广义动量映射（generalized moment map）}
定义广义动量映射 $\Phi_{\mathcal A}:M\to\mathcal A^*$。形如
$\Phi_{\mathcal A}^{-1}(p)$（$p\in\mathcal A^*$）的非空子集称为 $\mathcal A$ 的\emph{纤维}。
\end{definition}

与环面李群作用相伴的标准动量映射是这个定义的特例。因此，只要动量映射
$\Phi_{\mathcal A}$ 的纤维光滑，它便是\emph{余迷向}的。

\begin{theorem}[恩托夫和波尔捷罗维奇 (EnP06)]
\label{thm:22-nondisplaceable-fiber}
$C^\infty(M)$ 的任意有限维泊松可交换子空间至少有一个不可移纤维。此外，若每个
纤维都只有有限多个连通分支，则存在一个不可移的纤维连通分支。
\end{theorem}

\begin{proof}
反设 $\mathcal A$ 的所有纤维均可移。由于 $M$ 紧，$\Phi_{\mathcal A}$ 的像
$\Delta$ 紧。因而可选取 $\Delta$ 的有限开覆盖
$\mathcal U=\{U_1,\ldots,U_d\}$，使 $\Phi_{\mathcal A}^{-1}(U_i)$ 对
$i=1,\ldots,d$ 均可移。选取从属于 $\mathcal U$ 的单位分解
$\{\chi_i\}_{1\leq i\leq d}$。于是
$\Phi_{\mathcal A}^*\chi_i=\chi_i\circ\Phi_{\mathcal A}$ 的支集可移，故由消失性公理
$\zeta(\Phi_{\mathcal A}^*\chi_i)=0$。又有
$\{\Phi_{\mathcal A}^*\chi_i,\Phi_{\mathcal A}^*\chi_j\}=0$。应用归一化与部分可加性，
得到
\[
 1=\zeta(1)=\zeta\!\left(\sum_{i=1}^d\Phi_{\mathcal A}^*\chi_i\right)
 =\sum_{i=1}^d\zeta(\Phi_{\mathcal A}^*\chi_i)=0,
\]
矛盾。这证明了第一条陈述，证明完毕。

对第二条，只需进一步细化开集 $\Phi_{\mathcal A}^{-1}(U_i)$，分别覆盖纤维的各个连通
分支。细节留作习题。
\end{proof}

该定理也给出了获得具有辛交性质的余迷向对象的一种简单方法，因为只要
$\Phi_{\mathcal A}^{-1}(p)$ 光滑，它就是\emph{余迷向}子流形。

\begin{theorem}[恩托夫--波尔捷罗维奇]
\label{thm:22-nondisplaceable-skeleton}
考虑 $M$ 的一个三角剖分 $\Sigma$。若三角剖分 $\Sigma$ 的网格充分小，则 $M$ 的
$(2n-1)$-骨架 $\Sigma^{2n-1}$ 不可移。
\end{theorem}

\begin{proof}
先把 $\Sigma^{2n-1}$ 表成某个连续函数 $F$ 的水平集。如下定义 $F$：先在
$\Sigma^{2n-1}$ 上令 $F\equiv0$；然后在每个单形 $\Delta$ 上令
\[
 F=1-\operatorname{dist}(\,\cdot\,,x_\Delta)\geq0,
\]
其中 $x_\Delta$ 是 $\Delta$ 的重心，$\operatorname{dist}(\,\cdot\,,x_\Delta)$ 是
$\Delta$ 上从 $x_\Delta$ 出发的仿射距离函数。可选取 $\Delta$ 的分片线性坐标，使
所有重心都满足 $F(x_\Delta)=1$。依定义，$F^{-1}(0)=\Sigma^{2n-1}$。把 $\Sigma$
的网格取得充分小，可使 $\Sigma$ 中每个单形都可移。于是，对每个
$0<\epsilon\leq1$，$F^{-1}(\epsilon)$ 是“小”单形边界的不交并，故其每个连通分支
都包含在 $\Sigma$ 的某个单形中，从而可移。

现反设 $F^{-1}(0)=\Sigma^{2n-1}$ 可移，即存在 $\phi\in\Ham(M,\omega)$，使
\[
 \phi(F^{-1}(0))\cap F^{-1}(0)=\varnothing.
\]
\footnote{译者注：原书此式只印出
$\phi(F^{-1}(0))\cap F^{-1}(0)$，缺少表示“不相交”的 $=\varnothing$；译文按
“可移”的定义及下一句补出。}
由于 $F^{-1}(0)$ 紧，存在 $\epsilon_0>0$，使
\[
 \phi(F^{-1}(0))\subset F^{-1}((\epsilon_0,\infty)).
\]
由连续性，只要 $0<\epsilon_1<\epsilon_0$ 充分小，$F^{-1}([0,\epsilon_1])$ 就可移。
现用光滑函数 $G$ 在 $C^0$ 意义下逼近 $F$。若
$\lVert G-F\rVert_{C^0}$ 充分小，则存在
$\delta=\delta(\lVert G-F\rVert_{C^0})>0$ 及 $0<\epsilon_2<\epsilon_1$，使
\begin{enumerate}[label=(\arabic*)]
\item $\min G\geq-\delta$；
\item $G^{-1}([-\delta,\delta])\subset F^{-1}([0,\epsilon_0])$
      （特别地，$G^{-1}([-\delta,\delta])$ 可移）；
\item $G^{-1}((-\delta,\delta))\supset F^{-1}([0,\epsilon_2])$。
\end{enumerate}

考虑 $\epsilon>\delta$ 时的 $G^{-1}(\epsilon)$。若
$\lVert G-F\rVert_{C^0}$ 充分小，则 $F^{-1}(\epsilon)$ 与 $G^{-1}(\epsilon)$ 的
豪斯多夫距离可任意小，故 $G^{-1}(\epsilon)$ 的每个连通分支都必包含在
$\Sigma$ 的一个单形中，因而对所有 $\epsilon>\delta$ 都可移。另一方面，对任意
$\epsilon\leq\delta$，有
$G^{-1}(\epsilon)\subset G^{-1}([-\delta,\delta])$，故同样可移。显然，这与把
定理~\ref{thm:22-nondisplaceable-fiber} 应用于光滑函数 $G$ 的
$\mathcal A=\{G\}$ 相矛盾。
\end{proof}

本节余下部分用于证明定理~\ref{thm:22-spectral-partial-quasi-state}。为简化记号，仍记
\[
 \mathcal P_U^{\mathrm{ham}}
 =\mathcal P^{\mathrm{ham}}(\Symp_U(M,\omega),id).
\]
从 (22.3.25) 和\personBanyaga{} (Ba78) 的下述碎裂引理出发。

\begin{lemma}[碎裂引理 (Ba78)]
\label{lem:22-fragmentation}
令 $(M,\omega)$ 为闭辛流形。设 $\{U_1,\ldots,U_m\}$ 是一个开覆盖，其中每个
$U_i$ 上的 $\omega|_{U_i}$ 都正合。则任意
$\lambda\in\mathcal P^{\mathrm{ham}}(\Symp_0(M,\omega),id)$ 均可写成
\[
 \lambda=\prod_{i=1}^m\lambda_i,
 \qquad \lambda_i\in\mathcal P_{U_i}^{\mathrm{ham}},
 \qquad \operatorname{Cal}_{U_i}^{\mathrm{path}}(\lambda_i)=0.
\]
\end{lemma}

\begin{proof}
由显然的归纳论证，只需考虑 $N=2$，即 $U=U_1\cup U_2$。令道路
\[
 \lambda=\phi_H\in\mathcal P^{\mathrm{ham}}(\Symp_0(M,\omega),id).
\]
把单位区间分成
有限多个子区间并重新参数化到单位区间，可不失一般性地假设
$\lVert H\rVert_{C^1}<\epsilon$，其中充分小的正数 $\epsilon$ 只依赖于
$U_1,U_2,U$，其小量程度稍后确定。

取一对开子集 $V_1''\subset V_1'$，使
$V_1''\subset V_1'\subset U_1$、$\overline V_1''\subset V_1'\subset
\overline V_1'\subset U_1$，且 $V_1''\cup U_2\supset\operatorname{supp}H$。
令 $\chi:M\to[0,1]$ 为光滑截断函数，满足
$\operatorname{supp}\chi\subset U_1$ 且 $\chi=1$ 于 $V_1'$；考虑道路
$\lambda_1=\phi_{\chi H}$。容易看出，若 $\epsilon$ 充分小，则
$\lambda_1=\lambda$ 于 $V_1''$。此外，利用这些道路在 $C^1$ 意义下接近恒等道路，可
假设对所有 $t\in[0,1]$ 及 $x\notin V_1''\cup U_2$ 都有 $\lambda(t,x)\equiv x$。

特别地，$\lambda_2:=\lambda_1^{-1}\lambda$ 的支集包含在 $U_2$ 中，而
$\lambda_1$ 的支集包含在 $U_1$ 中。依定义 $\lambda=\lambda_1\lambda_2$，这便证明了
$m=2$ 的情形，从而完成证明。
\end{proof}

该引理引出如下定义。

\begin{definition}
\label{def:22-fragmentation-length}
给定开子集 $U\subset M$ 与 $\lambda\in\mathcal P_U^{\mathrm{ham}}$，定义
\stfhIndex{I090}{Suiliechangdu}{碎裂长度（fragmentation length）}
$\lambda$ 在 $U$ 上的\emph{碎裂长度}为把引理~\ref{lem:22-fragmentation} 应用于
$(U,\omega)$ 时乘积中因子的最小可能数目，记作 $\lVert\lambda\rVert_U$。
\end{definition}

下述命题正是 (EnP06) 定理 7.1 的道路版本。

\begin{proposition}
\label{prop:22-path-quasimorphism}
函数 $\mu_1^{\mathrm{path}}$（式~\eqref{eq:22-path-asymptotic-average}）良定，并满足：
\begin{enumerate}[label=\arabic*.]
\item （受控拟可加性）给定 $M$ 的可移开子集 $U$，存在只依赖于 $U$ 的常数 $K$，使
\begin{equation}
 \bigl|\mu_1^{\mathrm{path}}(\lambda\delta)
 -\mu_1^{\mathrm{path}}(\lambda)-\mu_1^{\mathrm{path}}(\delta)\bigr|
 \leq K\min\{\lVert\lambda\rVert_U,\lVert\delta\rVert_U\}
 \label{eq:22-controlled-quasi-additivity}
\end{equation}
对任意 $\lambda,\delta\in\mathcal P^{\mathrm{ham}}(\Symp(M,\omega),id)$ 成立；
\item （卡拉比性质）对任意 $\lambda\in\mathcal P_U^{\mathrm{ham}}$，
$\mu_1^{\mathrm{path}}(\lambda)=\operatorname{Cal}_U^{\mathrm{path}}(\lambda)$。
\end{enumerate}
\end{proposition}

\begin{proof}
由命题 22.3.3，对任意 $\lambda\in\mathcal P_U^{\mathrm{ham}}$ 有
\begin{equation}
 \rho(\lambda;1)+\rho(\lambda^{-1};1)\leq2e_\rho(U;M).
 \label{eq:22-path-inverse-bound}
\end{equation}
记 $C=2e_\rho(U;M)$。令 $\delta$ 为 $M$ 上任意哈密顿道路。由三角不等式，
\begin{align}
 \rho(\lambda\delta;1)
 &\geq\rho(\delta;1)-\rho(\lambda^{-1};1)\notag\\
 &\geq\rho(\delta;1)-(C-\rho(\lambda;1))\notag\\
 &\geq\rho(\delta;1)+\rho(\lambda;1)-C.
 \label{eq:22-path-lower-triangle}
\end{align}
三角不等式还给出
$\rho(\lambda\delta;1)\leq\rho(\lambda;1)+\rho(\delta;1)$。

从~\eqref{eq:22-path-lower-triangle} 出发，对
$\lambda_1,\ldots,\lambda_m\in\mathcal P_U^{\mathrm{ham}}$ 作归纳，得到
\begin{equation}
 \left|\rho(\lambda_1\cdots\lambda_m\delta;1)
 -\sum_{i=1}^m\rho(\lambda_i;1)-\rho(\delta;1)\right|\leq mC.
 \label{eq:22-fragment-product-bound}
\end{equation}
特别地，取 $\delta\equiv id$ 并取第 $\ell$ 次幂，有
\begin{equation}
 \left|\rho((\lambda_1\cdots\lambda_m)^\ell;1)
 -\ell\sum_{i=1}^m\rho(\lambda_i;1)\right|\leq\ell mC.
 \label{eq:22-fragment-power-bound}
\end{equation}

现任取哈密顿道路 $\lambda$，由碎裂引理把它写成
$\lambda=\lambda_1\cdots\lambda_m$，且每个 $i$ 都有
$\lVert\lambda_i\rVert=1$。于是
\[
 \rho(\lambda^\ell;1)
 \geq\ell\sum_{i=1}^m\rho(\lambda_i;1)-\ell mC
 =\ell\left(\sum_{i=1}^m\rho(\lambda_i;1)-mC\right),
\]
从而
\[
 \frac{\rho(\lambda^\ell;1)}{\ell}
 \geq\sum_{i=1}^m\rho(\lambda_i;1)-mC=:E.
\]
$m$ 及该分解只依赖于 $\lambda$，故不等式右端只依赖于 $\lambda$。因此对所有
$\ell$，
\begin{equation}
 E\leq\frac{\rho(\lambda^\ell;1)}{\ell}\leq\rho(\lambda;1),
 \label{eq:22-asymptotic-bounds}
\end{equation}
其中第二个不等式来自三角不等式。下述引理将完成极限良定性的证明，从而证明
$\mu_1^{\mathrm{path}}(\lambda)$ 良定。

\begin{lemma}
\label{lem:22-asymptotic-limit}
当 $\ell\to\infty$ 时，极限
\[
 \lim_{\ell\to\infty}\frac{\rho(\lambda^\ell;1)}{\ell}
\]
存在。
\end{lemma}

\begin{proof}
仿照 (C09) 引理 2.21 的证明。由~\eqref{eq:22-asymptotic-bounds}，
\[
 0\leq\rho(\lambda^\ell;1)-E\ell.
\]
把 $\ell$ 取为正整数 $i$ 的二进数 $\ell=2^i$，再应用三角不等式，得到
\[
 0\leq\rho(\lambda^{2^i};1)-2^iE
 \leq2\rho(\lambda^{2^{i-1}};1)-2^iE.
\]
除以 $2^i$，得到
\[
 0\leq\frac{\rho(\lambda^{2^i};1)-2^iE}{2^i}
 \leq\frac{\rho(\lambda^{2^{i-1}};1)-2^{i-1}E}{2^{i-1}}.
\]
所以序列 $(\rho(\lambda^{2^i};1)-2^iE)/2^i$ 非负，并随 $i\to\infty$ 单调递减，
因而收敛。显然，这蕴含 $i\to\infty$ 时 $\rho(\lambda^{2^i};1)/2^i$ 收敛，继而
蕴含 $\ell\to\infty$ 时 $\rho(\lambda^\ell;1)/\ell$ 收敛。
\end{proof}

接着证明受控拟可加性~\eqref{eq:22-controlled-quasi-additivity}。考虑不为恒等道路的任意
$\lambda,\delta$，其中
\[
 \lambda,\delta\in\mathcal P^{\mathrm{ham}}(\Symp_0(M,\omega),id).
\]
此时
$\lVert\lambda\rVert_U,\lVert\delta\rVert_U\geq1$。我们断言
\begin{equation}
 \bigl|\mu_1^{\mathrm{path}}(\lambda\delta)
 -\mu_1^{\mathrm{path}}(\lambda)-\mu_1^{\mathrm{path}}(\delta)\bigr|
 \leq2C\min\{2\lVert\lambda\rVert_U-1,2\lVert\delta\rVert_U-1\}.
 \label{eq:22-controlled-induction}
\end{equation}
\footnote{译者注：原书式 (22.4.40) 的右端印作
$2C(2\lVert\lambda\rVert_U-1,2\lVert\delta\rVert_U-1)$，两个量之间只有逗号，
没有运算符。下文明确以
$m=\min\{\lVert\lambda\rVert_U,\lVert\delta\rVert_U\}$ 归纳，故译文补出
$\min$。}
对 $m:=\min\{\lVert\lambda\rVert_U,\lVert\delta\rVert_U\}$ 作归纳。$m=0$，即
$\lambda$ 或 $\delta$ 为恒等道路时，结论平凡。因此从 $m=1$ 开始；即
$\lVert\lambda\rVert_U=1$ 或 $\lVert\delta\rVert_U=1$。不失一般性，设
$\lVert\lambda\rVert_U=1$。写成
\begin{equation}
 (\lambda\delta)^k
 =\left(\prod_{i=0}^{k-1}\delta^i\lambda\delta^{-i}\right)\delta^k.
 \label{eq:22-conjugate-product}
\end{equation}
由 $\lVert\,\cdot\,\rVert_U$ 的共轭不变性，
$\lVert\delta^i\lambda\delta^{-i}\rVert_U=\lVert\lambda\rVert_U=1$。应用
\eqref{eq:22-fragment-product-bound} 及 $\rho(\,\cdot\,;1)$ 的共轭不变性，得到
\[
 |\rho((\lambda\delta)^k;1)-k\rho(\lambda;1)-\rho(\delta^k;1)|\leq Ck.
\]
将它与不等式
\[
 |\rho(\lambda^k;1)-k\rho(\lambda;1)|\leq Ck
\]
结合；原书称后者对应于式 (22.4.36) 的 $m=1,\ell=k$ 情形。
\footnote{译者注：这里的 $m,\ell$ 只出现在式 (22.4.38)，而式 (22.4.36) 并无这两个
参数；原书交叉引用疑应为 (22.4.38)。译文保留原指称并说明疑点。}
于是
\[
 \frac1k\left|
 (\rho((\lambda\delta)^k;1)-k\rho(\lambda;1)-\rho(\delta^k;1))
 -(\rho(\lambda^k;1)-k\rho(\lambda;1))\right|\leq2C,
\]
即
\[
 \frac1k\left|\bigl(\rho((\lambda\delta)^k;1)-\rho(\delta^k;1)\bigr)
 -\rho(\lambda^k;1)\right|\leq2C.
\]
令 $k\to\infty$，得到
\[
 |\mu_1^{\mathrm{path}}(\lambda\delta)-\mu_1^{\mathrm{path}}(\lambda)
 -\mu_1^{\mathrm{path}}(\delta)|\leq2C,
\]
这就以 $K=2C$ 完成 $m=1$ 的证明。

现设~\eqref{eq:22-controlled-induction} 对所有 $m\leq N$ 成立，并考虑 $m=N+1$。
仍可设 $\lVert\lambda\rVert_U=N+1$。把 $\lambda$ 分解为
$\lambda=\lambda'\lambda_1$，其中 $\lVert\lambda_1\rVert_U=1$、
$\lVert\lambda'\rVert_U=N$。由归纳假设，
\begin{align*}
 |\mu_1^{\mathrm{path}}(\lambda'\lambda_1\delta)
 -\mu_1^{\mathrm{path}}(\lambda')-\mu_1^{\mathrm{path}}(\lambda_1\delta)|
 &\leq2C(2m-1),\\
 |\mu_1^{\mathrm{path}}(\lambda_1\delta)-\mu_1^{\mathrm{path}}(\lambda_1)
 -\mu_1^{\mathrm{path}}(\delta)|&\leq2C,\\
 |\mu_1^{\mathrm{path}}(\lambda_1)+\mu_1^{\mathrm{path}}(\lambda')
 -\mu_1^{\mathrm{path}}(\lambda'\lambda_1)|&\leq2C.
\end{align*}
将这些不等式相加，得到
\[
 |\mu_1^{\mathrm{path}}(\lambda\delta)-\mu_1^{\mathrm{path}}(\lambda)
 -\mu_1^{\mathrm{path}}(\delta)|\leq2C(2m+1).
\]
取 $K=4C=8e_\rho(U;M)$，便完成受控拟可加性的证明。

最后证明 $\mu_1^{\mathrm{path}}$ 的卡拉比性质。令
$\lambda=\phi_F\in\mathcal P_U^{\mathrm{ham}}$，并选取哈密顿道路
$\delta=\phi_H$，使 $\phi_H^1(U)\cap\overline U=\varnothing$。命题 22.3.3 给出
\[
 \rho(\phi_H\phi_F^m;1)=\rho(\phi_H;1)
 +\operatorname{Cal}_U^{\mathrm{path}}(\phi_F^m)
 =\rho(\phi_H;1)+m\operatorname{Cal}_U^{\mathrm{path}}(\phi_F)
\]
对所有 $m$ 成立。另一方面，由~\eqref{eq:22-path-lower-triangle} 还有
\footnote{译者注：原书在此处第一项卡拉比不变量中印作
$\operatorname{Cal}_U^{\mathrm{path}}(\phi_F^m;1)$；卡拉比不变量并无第二个参数，且
紧接着的等式也按单参数使用，故译文删去多出的 ``$;1$''。}
\[
 |\rho(\phi_H\phi_F^m;1)-\rho(\phi_H;1)-\rho(\phi_F^m;1)|
 \leq2e_\rho(U;M).
\]
除以 $m$ 并令 $m\to\infty$，得到
\begin{align*}
 \lim_{m\to\infty}\frac{\rho(\phi_F^m;1)}m
 &=\lim_{m\to\infty}\frac{\rho(\phi_H\phi_F^m;1)}m\\
 &=\lim_{m\to\infty}\frac{\rho(\phi_H;1)}m
 +\operatorname{Cal}_U^{\mathrm{path}}(\phi_F)
 =\operatorname{Cal}_U^{\mathrm{path}}(\phi_F).
\end{align*}
因此由 $\mu_1^{\mathrm{path}}$ 的定义，
\[
 \mu_1^{\mathrm{path}}(\phi_F)
 =\lim_{m\to\infty}\frac{\rho(\phi_F^m;1)}m
 =\operatorname{Cal}_U^{\mathrm{path}}(\phi_F).
\]
这完成证明。
\end{proof}

现在可以完成定理~\ref{thm:22-spectral-partial-quasi-state} 的证明。

单调性、归一化、利普希茨连续性及辛不变性直接来自
$\rho(\,\cdot\,;1)$ 的公理，其证明留给读者。半齐次性直接来自
$\mu_1^{\mathrm{path}}$ 的齐次性与利普希茨连续性。

只需再证部分可加性与消失性。令 $F,H\in C^\infty(M)$ 满足
$\{F,H\}=0$ 且 $\operatorname{supp}H\subset U$。
\footnote{译者注：原书此处印作 $\operatorname{supp}H\in U$；支集是子集而非点，故译文
依本段语义及卡拉比性质的适用条件改为 $\operatorname{supp}H\subset U$。}
注意对所有 $k\in\mathbb N$ 都有
$\lVert\phi_H^k\rVert_U=1$。由于 $\{F,H\}=0$，流 $\phi_H$ 与 $\phi_F$ 可交换，故
\[
 \mu_1^{\mathrm{path}}(\phi_F\phi_H)
 =\frac1k\mu_1^{\mathrm{path}}((\phi_F\phi_H)^k)
 =\frac1k\mu_1^{\mathrm{path}}(\phi_F^k\phi_H^k).
\]
而由~\eqref{eq:22-controlled-quasi-additivity}，对所有 $k$ 有
\[
 |\mu_1^{\mathrm{path}}(\phi_F^k\phi_H^k)
 -\mu_1^{\mathrm{path}}(\phi_F^k)-\mu_1^{\mathrm{path}}(\phi_H^k)|\leq K.
\]
结合两式并除以 $k$，得到
\[
 |\mu_1^{\mathrm{path}}(\phi_F\phi_H)-\mu_1^{\mathrm{path}}(\phi_F)
 -\mu_1^{\mathrm{path}}(\phi_H)|\leq\frac Kk.
\]
这对所有 $k$ 成立，因而
\[
 \mu_1^{\mathrm{path}}(\phi_F\phi_H)
 =\mu_1^{\mathrm{path}}(\phi_F)+\mu_1^{\mathrm{path}}(\phi_H).
\]
由卡拉比性质，
\[
 \mu_1^{\mathrm{path}}(\phi_H)
 =\operatorname{Cal}_U^{\mathrm{path}}(\phi_H)
 =\operatorname{vol}_\omega(M)^{-1}\int_M H\,d\mu_\omega.
\]
把它代入 $\zeta_{\mathrm{sp}}$ 的定义，得到
$\zeta_{\mathrm{sp}}(F+H)=\zeta_{\mathrm{sp}}(F)$。

\section{恩托夫和波尔捷罗维奇的卡拉比拟同态}

由\personBanyaga{}定理，闭 $M$ 的群 $\Ham(M,\omega)$ 是单群，因而不存在从
$\Ham(M,\omega)$ 到 $\mathbb R$（或任意非平凡交换群）的非平凡满同态。换言之，
对任意\emph{满射} $\mu:\Ham(M,\omega)\to\mathbb R$，缺陷
\[
 \sup_{\phi_1,\phi_2\in\Ham(M,\omega)}
 |\mu(\phi_1\phi_2)-\mu(\phi_1)-\mu(\phi_2)|
\]
为正。若 $\mu$ 是同态，则其缺陷当然为零。当不存在非平凡满同态时，退而求其次，
可以期待缺陷\emph{有界}的映射 $\mu$。事实表明，$\Ham(M,\omega)$ 带有一些源自哈密顿流及其
不动点动力学的有趣拟同态。恩托夫和波尔捷罗维奇 (EnP03) 以谱不变量的公理性质为
基础完成了这一构造。这再次体现出：谱不变量编码 $(M,\omega)$ 的辛拓扑性质以及
$(M,\omega)$ 上哈密顿流及其周期轨道的动力学性质，并解码蕴藏在弗洛尔同调解析理论
中的辛不变量。

在恩托夫和波尔捷罗维奇的构造中，关键是对适当选取的\emph{幂等元}
$e\in QH^*(M)$ 证明赋值 $\phi_H\mapsto\rho(\phi_H;e)$ 为拟同态。他们在底层辛流形
$(M,\omega)$ 的量子上同调环满足某种代数性质（例如半单性），且弗罗贝尼乌斯配对定义
非退化配对的假设下证明了这一拟同态性质。后一个非退化性对\emph{有理}辛流形不难证明，但
对无理情形困难得多（见 (Ush10b)）。量子上同调环的半单性对 $S^2$ 成立，更一般地
对 $\mathbb{CP}^n$ 成立。这把量子上同调环的代数性质与借助谱不变量公理研究哈密顿
动力学卓越地融合在一起。即使对 $S^2$，它也给出了保面积微分同胚群上拟同态的新构造；
这类研究在动力系统领域受到广泛关注，例如见 (GG97), (GG04), (Gh06), (C09)。

\subsection{再论庞加莱对偶}

先汇集弗洛尔同调中对偶性的若干基本事实。考虑滤过弗洛尔复形
\[
 CF_k^\lambda(H)=CF_k^{(-\infty,\lambda]}(H)
 =\{\alpha\in CF_k(H)\mid\lambda_H(\alpha)\leq\lambda\}
\]
及其商
\[
 CF_k^{(\mu,\lambda]}(H)
 =CF_k^{(-\infty,\lambda]}(H)/CF_k^{(-\infty,\mu]}(H),
\]
其中 $H$ 为任意一般的非退化哈密顿量，$k\in\mathbb Z_+$。以
\[
 i_\lambda:CF_*^{(-\infty,\lambda]}(H)\longrightarrow CF_*(H)
\]
表示典范包含，以
\[
 \pi_\lambda:CF_*(H)\longrightarrow CF_*^{(\lambda,\infty)}(H)
 =CF_*(H)/CF_*^{(-\infty,\lambda]}(H)
\]
表示典范投影。对每个固定的 $\lambda_1,\lambda_2<\lambda$，存在自然投影
\[
 \pi^\lambda_{\lambda_2\lambda_1}:CF_k^{(\lambda_1,\lambda]}(H)
 \longrightarrow CF_k^{(\lambda_2,\lambda]}(H),
 \qquad \lambda_1\leq\lambda_2,
\]
它们定义一个逆系统。于是有下述引理。

\begin{lemma}
\label{lem:22-filtered-inverse-limit}
令 $(M,\omega)$ 为任意紧辛流形，$(H,J)$ 弗洛尔正则。则典范极限映射
\[
 \pi^\lambda:CF_k^{(-\infty,\lambda]}(H)
 \longrightarrow\varprojlim_{\mu}CF_k^{(\mu,\lambda]}(H)
\]
对所有 $\lambda\in\mathbb R\cup\{+\infty\}$ 都是同构。
\end{lemma}

\begin{proof}
令 $\alpha\in CF_k^{(-\infty,\lambda]}(H)$。依定义，可写成
\[
 \alpha=\sum_{i=1}^\infty a_i[z_i,w_i]
 =\lim_{N\to\infty}\sum_{i=1}^N a_i[z_i,w_i],
 \qquad
 \mathcal A_H([z_i,w_i])\geq\mathcal A_H([z_{i+1},w_{i+1}]),
\]
其中或者只有有限多个 $a_i$ 非零，或者
$\lim_{i\to\infty}\mathcal A_H([z_i,w_i])=-\infty$。由此，$\alpha=0$ 当且仅当
$\sum_{i=1}^N a_i[z_i,w_i]=0$ 对所有 $N$ 成立，也即
$\pi^\lambda_{-\infty\mu}(\alpha)=0$ 对所有 $\mu$ 成立。因此 $\pi^\lambda$ 单射。

另一方面，$\varprojlim_\mu CF_k^{(\mu,\lambda]}(H)$ 中任意元素
$\{\alpha_\mu\}$ 都可由有限和
\[
 \sum_{i=1}^{N_\mu}a_i[z_i,w_i]
\]
表示，其中 $N_\mu\in\mathbb Z_+$。
\footnote{译者注：原书证明此处把逆极限中的区间印作 $[\mu,\lambda)$，与本引理陈述、
本小节此前定义的 $(\mu,\lambda]$ 以及后续投影均不一致；译文按已定义的
$(\mu,\lambda]$ 统一。}
设给定满足
$\pi^\lambda_{\mu'\mu}(\alpha_{\mu'})=\alpha_\mu$ 的元素
$\{\alpha_\mu\}_{\mu\in\mathbb R}$。对任意满足 $k\to\infty$ 时
$\mu_k\to-\infty$ 的序列，可写
\[
 \alpha_{\mu_k}=\sum_{i=1}^{N_{\mu_k}}a_i[z_i,w_i],
\]
它收敛到 $CF^\lambda(H)$ 中某个 $\alpha$。容易验证该极限与序列 $\mu_k$ 的选择无关，
故 $\pi^\lambda$ 为同构。证明完毕。
\end{proof}


\begin{remark}
\label{rem:22-rational-finite-window}
现设 $(M,\omega)$ 有理。此时 $CF_k^{(\mu,\lambda]}(H)$ 是有限维
$\mathbb Q$-向量空间。需要强调，对无理的 $(M,\omega)$ 并非如此。
\end{remark}

考虑 $H$ 与 $\widetilde H=-H(1-t,x)$，并定义 $\mathbb Q$-值配对
\begin{equation}
 \langle\,\cdot\,,\,\cdot\,\rangle:
 CF_n^{(-\infty,\lambda]}(H)\times
 CF_n^{(-\lambda,\infty)}(\widetilde H)\longrightarrow\mathbb Q
 \label{eq:22-filtered-pairing}
\end{equation}
如下：
\[
 \langle\alpha,\widetilde\beta\rangle
 =\sum_{[z,w]\in\operatorname{Crit}\mathcal A_H}
 a_{[z,w]}b_{[\widetilde z,\widetilde w]},
\]
其中
$\alpha=\sum_{[z,w]\in\operatorname{Crit}\mathcal A_H}a_{[z,w]}[z,w]$，而
$\widetilde\beta=
\sum_{[\widetilde z,\widetilde w]\in\operatorname{Crit}\mathcal A_{\widetilde H}}
b_{[\widetilde z,\widetilde w]}[\widetilde z,\widetilde w]$。由诺维科夫有限性条件，
\[
 CF_n^{(-\lambda,\infty)}(\widetilde H)
 =CF_n(\widetilde H)/CF_n^{(-\infty,-\lambda]}(\widetilde H)
\]
中的任意元素都可用有限和表示，故上述配对良定且非退化。它进而诱导单射链同态
\begin{equation}
 (CF_n^{(-\infty,\lambda]}(H),\partial_{(H,J)})
 \longrightarrow
 \bigl(\operatorname{Hom}_{\mathbb Q}(CF_n^{(-\lambda,\infty)}(\widetilde H),
 \mathbb Q),(\partial_{(\widetilde H,\widetilde J)})^*\bigr).
 \label{eq:22-filtered-chain-duality-map}
\end{equation}
回忆典范等同
\[
 \operatorname{Hom}(CF_n^{(-\lambda,\infty)}(\widetilde H),\mathbb Q)
 =\{\varphi\in\operatorname{Hom}(CF_n(\widetilde H),\mathbb Q)
 \mid \varphi|_{CF_n^{(-\infty,-\lambda]}(\widetilde H)}\equiv0\}.
\]
\footnote{译者注：原书此式限制下标印作
$CF_n^{(-\infty,\lambda]}(\widetilde H)$，但左端是商空间
$CF_n(\widetilde H)/CF_n^{(-\infty,-\lambda]}(\widetilde H)$ 的对偶；泛函应在后一个
子空间上消失，故译文补上遗漏的负号。}
记
$\operatorname{Hom}(CF_n^{(-\lambda,\infty)}(\widetilde H),\mathbb Q)
=CF_n^{(-\lambda,\infty)}(\widetilde H)^*$。通过这一等同，当 $\lambda>\mu$ 时有典范
包含映射
\[
 CF_n^{(-\mu,\infty)}(\widetilde H)^*\lhook\joinrel\longrightarrow
 CF_n^{(-\lambda,\infty)}(\widetilde H)^*.
\]
此外还有自然等同
\[
 CF_n^{(-\lambda,-\mu]}(\widetilde H)^*
 \cong CF_n^{(-\lambda,\infty)}(\widetilde H)^*/
 CF_n^{(-\mu,\infty)}(\widetilde H)^*.
\]

下述定理是恩托夫和波尔捷罗维奇对有理 $(M,\omega)$ 构造拟同态时所用的关键事实。
无理情形的相应定理难证得多；厄舍 (Ush10b) 运用非阿基米德分析的论证证明了它。
这里按照恩托夫和波尔捷罗维奇给出有理情形证明的全部细节；无理情形请读者参阅
厄舍的原文。

\begin{theorem}
\label{thm:22-rational-filtered-duality}
令 $(M,\omega)$ 为任意紧有理辛流形。则典范 $\mathbb Q$-同态
\eqref{eq:22-filtered-chain-duality-map} 是链同构。特别地，配对
\eqref{eq:22-filtered-pairing} 下降为非退化配对
\begin{equation}
 HF_n^{(-\infty,\lambda]}(H)\times
 HF_n^{(-\lambda,\infty)}(\widetilde H)\longrightarrow\mathbb Q.
 \label{eq:22-filtered-homology-pairing}
\end{equation}
\end{theorem}

\begin{proof}
由注~\ref{rem:22-rational-finite-window}，
$CF_k^{(\mu,\lambda]}(H)$ 与
$CF_n^{(-\lambda,-\mu]}(\widetilde H)^*$ 都是有限维 $\mathbb Q$-向量空间（在无理
情形下它们为无限维，这正是本证明不能覆盖无理情形的主要原因）。有链复形交换图
\begin{equation}
\begin{array}{ccc}
 CF_*^{(-\infty,\lambda]}(H)&\xrightarrow{\ \pi_\mu^\lambda\ }&
 CF_*^{(\mu,\lambda]}(H)\\[3pt]
 \big\downarrow&&\big\downarrow\\[3pt]
 CF_*^{(-\lambda,\infty)}(\widetilde H)^*&\longrightarrow&
 CF_*^{(-\lambda,-\mu]}(\widetilde H)^* .
\end{array}
 \label{eq:22-filtered-chain-diagram}
\end{equation}
\footnote{译者注：原书在此交换图右下角以及下文所取同调的右侧复形中把链次数印作
固定的 $n$，但上下文讨论的是链复形，紧接着的文字也明确写成 $CF_*$。译文在这些
复形位置统一改为 $*$；定理陈述中固定次数 $n$ 的配对保持不变。}
映射
$CF_*^{(\mu,\lambda]}(H)\to CF_*^{(-\lambda,-\mu]}(\widetilde H)^*$ 是链同构：
在每个次数上该映射单射，而由 $(M,\omega)$ 的\emph{有理性}，两个空间均有限维。因此上图
诱导交换图
\[
\begin{array}{ccc}
 CF_*^{(-\infty,\lambda]}(H)&\xrightarrow{\ \pi^\lambda\ }&
 \varprojlim_\mu CF_*^{(\mu,\lambda]}(H)\\[3pt]
 \big\downarrow&&\big\downarrow\\[3pt]
 CF_*^{(-\lambda,\infty)}(\widetilde H)^*&\longrightarrow&
 \varprojlim_\mu CF_*^{(-\lambda,-\mu]}(\widetilde H)^* .
\end{array}
\]
其中右侧向下映射是链同构。水平映射均为同构，故左侧向下映射也是链同构。因此它
诱导同构
\[
 H\bigl(CF_*^{(-\infty,\lambda]}(H),\partial_{(H,J)}\bigr)
 \cong
 H\bigl(CF_*^{(-\lambda,\infty)}(\widetilde H)^*,
 \partial_{(\widetilde H,\widetilde J)}^*\bigr).
\]
左端恰为 $HF^{(-\infty,\lambda]}(H)$；由域 $\mathbb Q$ 上的泛系数定理，右端同构于
\[
 \operatorname{Hom}_{\mathbb Q}
 \bigl(H(CF_*^{(-\lambda,\infty)}(\widetilde H),
 \partial_{(\widetilde H,\widetilde J)}),\mathbb Q\bigr)
 =HF^{(-\lambda,\infty)}(\widetilde H,\widetilde J).
\]
这一同构自然诱导对所有 $\lambda$ 都成立的（强）非退化配对
\[
 HF_n^{(-\infty,\lambda]}(H)\times
 HF_n^{(-\lambda,\infty)}(\widetilde H)\longrightarrow\mathbb Q,
\]
故得到~\eqref{eq:22-filtered-homology-pairing}。
\end{proof}

当 $(M,\omega)$ 无理时，上述映射
\[
 CF_*^{(\mu,\lambda]}(H)\longrightarrow
 CF_*^{(-\lambda,-\mu]}(\widetilde H)^*
\]
不是同构。事实上，后一个空间大得多：该映射单射但不满射。

\begin{example}
\label{ex:22-irrational-duality-failure}
设 $(M,\omega)$ 无理，使周期群
$\{\omega(A)\mid A\in\pi_2(M)\}$ 是 $\mathbb R$ 的稠密子群。则对任意哈密顿量
$H$，$\operatorname{Spec}(H)$ 是 $\mathbb R$ 的稠密子集。考虑形式和
\[
 \gamma=\sum_{i=1}^\infty a_i[z_i,w_i],\qquad a_i\neq0,
\]
满足所有 $i$ 都有 $\mathcal A_H([z_i,w_i])\in(\mu,\lambda)$，且当 $i\neq j$ 时
$[z_i,w_i]\neq[z_j,w_j]$。由于该无穷和不满足诺维科夫有限性条件，
$\gamma\notin CF_*^{(\mu,\lambda]}(H)$。但容易验证，配对
\[
 \gamma\longmapsto\langle\gamma,\,\cdot\,\rangle
\]
定义了 $CF_*^{(-\lambda,-\mu)}(\widetilde H)$ 上不属于
$CF_*^{(\mu,\lambda]}(H)$ 之像的线性泛函。
\end{example}

现在陈述 $\widetilde H$ 与 $H$ 的谱不变量之间关系的一般定理。该关系在恩托夫和
波尔捷罗维奇的拟同态构造中起关键作用。

\begin{lemma}[(EnP03) 的引理 2.2]
\label{lem:22-spectral-duality}
令 $\Pi$ 为式 (20.4.43) 给出的配对。对
$b\in QH^*(M)\setminus\{0\}$，以 $\Upsilon(b)$ 表示所有满足
$\Pi(b,a)\neq0$ 的 $a\in QH^*(M)$ 所成的集合。则
\begin{equation}
 \rho(\widetilde H;b)=-\inf_{a\in\Upsilon(b)}\rho(H;a).
 \label{eq:22-spectral-duality}
\end{equation}
\end{lemma}

\begin{proof}
先证
\begin{equation}
 \rho(\widetilde H;b)\geq-\inf_{a\in\Upsilon(b)}\rho(H;a).
 \label{eq:22-spectral-duality-lower}
\end{equation}
$\rho$ 的\emph{三角不等式}是证明该不等式的关键。

反设
$\rho(\widetilde H;b)+\inf_{a\in\Upsilon(b)}\rho(H;a)<0$。于是存在满足
$\Pi(a,b)\neq0$ 的 $a$，使
\[
 \rho(\widetilde H;b)+\rho(H;a)<0.
\]
另一方面，对任意满足 $\Pi(a,b)\neq0$ 的 $a$，由 $QH^*(M)$ 的弗罗贝尼乌斯性质
(20.4.43)，有 $(a*b)_0\neq0$，从而
$\rho(0;a*b)=-\nu(a*b)\geq0$。但这给出
\[
 0\leq\rho(0;a*b)=\rho(H\#\widetilde H;a*b)
 \leq\rho(H;a)+\rho(\widetilde H;b)<0,
\]
矛盾，故~\eqref{eq:22-spectral-duality-lower} 成立。

接着证明反向不等式，即
\begin{equation}
 \inf_{a\in\Upsilon(b)}\rho(H;a)\leq-\rho(\widetilde H;b).
 \label{eq:22-spectral-duality-upper}
\end{equation}
上面的\emph{对偶配对非退化性}在证明中起关键作用。

令 $c=\rho(\widetilde H;b)>-\infty$，并令
$\Phi_{\widetilde H}:QH^*(M)\to FH_*(\widetilde H)$ 为 PSS 映射。对任意
$\epsilon>0$，有
$\Phi_{\widetilde H}(b)\notin\operatorname{Im}(i_{c-\epsilon})_*
\subset HF_*(\widetilde H)$。由序列
\[
 HF_*^{c-\epsilon}(\widetilde H)
 \xrightarrow{(i_{c-\epsilon})_*}HF_*(\widetilde H)
 \xrightarrow{(\pi_{c-\epsilon})_*}HF_*^{(c-\epsilon,\infty)}(\widetilde H)
\]
在中间项的正合性，
$\ker(\pi_{c-\epsilon})_*=\operatorname{Im}(i_{c-\epsilon})_*$，故
\[
 (\pi_{c-\epsilon})_*(\Phi_{\widetilde H}(b))\neq0.
\]
于是由对偶配对的非退化性，可找到弗洛尔同调类
$0\neq a_H\in HF_*^{(-\infty,-c+\epsilon]}(H)$，它由闭链
$\alpha_H\in CF_*(H)$ 表示，满足 $\lambda_H(\alpha_H)\leq-c+\epsilon$，且
\[
 L'([\alpha_H],\Phi_{\widetilde H}(b))\neq0.
\]
把弗洛尔同调类 $[\alpha_H]\in FH_*(H)$ 写成
$[\alpha_H]=\Phi_H(a')$，其中 $a'\in QH^*(M)$。于是
\[
 \Pi(a',b)=L'(\Phi_H(a'),\Phi_{\widetilde H}(b))
 =L'([\alpha_H],\Phi_{\widetilde H}(b))\neq0,
\]
故 $a'\in\Upsilon(b)$。最后一个等号使用了定理~\ref{thm:22-rational-filtered-duality}
与~\eqref{eq:22-filtered-pairing}。

由于 $\rho(H;a')\leq\lambda_H(\alpha_H)$，有
$\rho(H;a')\leq-c+\epsilon$，从而
\[
 \inf_{a\in\Upsilon(b)}\rho(H;a)\leq\rho(H;a')\leq-c+\epsilon.
\]
这对所有 $\epsilon>0$ 成立，故~\eqref{eq:22-spectral-duality-upper} 成立，引理得证。
\end{proof}

\subsection{拟同态的一般性质}

先概述一般拟同态的基本事实。\personCalegari{}的著作 (C09) 是一般语境下拟同态研究的一本
较新而出色的参考书。考虑任意群 $G$ 与满射 $\mu:G\to\mathbb R$。

\begin{definition}[拟同态]
\label{def:22-quasimorphism}
\stfhSubIndex{I163}{Nitongtai}{拟同态（quasimorphism）}{Quexian}{缺陷（defect）}
给定函数 $\mu:G\to\mathbb R$，定义其（同态）缺陷
$\operatorname{Def}_\mu:G\times G\to\mathbb R$ 为
\[
 \operatorname{Def}_\mu(g_1,g_2)=\mu(g_1g_2)-\mu(g_1)-\mu(g_2),
\]
其范数为
\[
 D(\mu):=|\operatorname{Def}_\mu|
 =\sup_{g_1,g_2\in G}|\mu(g_1g_2)-\mu(g_1)-\mu(g_2)|.
\]
若 $D(\mu)<\infty$，即存在 $R>0$，使
\begin{equation}
 |\mu(g_1g_2)-\mu(g_1)-\mu(g_2)|<R,
 \label{eq:22-quasimorphism-defect}
\end{equation}
则称 $\mu$ 为拟同态。
\end{definition}

映射 $\operatorname{Def}$ 衡量拟同态 $\mu:G\to\mathbb R$ 的同态性质之缺陷。
回忆群上同调的如下标准定义（例如见 (Wb94)）。

\begin{definition}
\label{def:22-group-cohomology}
令 $M$ 为左 $G$-模。设 $(C^*(G,M),\delta)$ 为复形，其中 $C^n(G,M)$ 是函数
$\varphi:G^n\to M$ 的模，而上边界映射
$\delta:C^n(G,M)\to C^{n+1}(G,M)$ 由
\begin{align}
 (\delta\varphi)(g_1,\ldots,g_{n+1})
 &=g_1\cdot\varphi(g_2,\ldots,g_{n+1})\notag\\
 &\quad+\sum_{i=1}^n(-1)^i
 \varphi(g_1,\ldots,g_{i-1},g_ig_{i+1},g_{i+2},\ldots,g_{n+1})\notag\\
 &\quad+(-1)^{n+1}\varphi(g_1,\ldots,g_n)
 \label{eq:22-group-coboundary}
\end{align}
给出。若函数 $\varphi:G^k\to M$ 满足 $\delta\varphi=0$，则称其为群 $k$-余闭链；
若对某个余链 $\psi:G^{k-1}\to M$ 有 $\varphi=\delta\psi$，则称其为余边界。
以左 $G$-模 $M$ 为系数的 $G$ 的上同调群，是复形 $(C^*(G,M),\delta)$ 的上同调群，
记作 $H^*(G,M)$。
\footnote{译者注：原书此定义对一般左 $G$-模 $M$ 却把余链值域两次印作
$\mathbb R$，并把“上同调群”称为该复形的“同调群”；译文按标准群上同调定义及式中
$g_1$ 对值的 $M$-作用，改为值域 $M$ 与“上同调群”。}
\end{definition}

考虑 $M=\mathbb R$ 且 $G$ 作用平凡的特例。此时~\eqref{eq:22-group-coboundary} 化为
\[
 (\delta\varphi)(g_1,g_2,g_3)
 =\varphi(g_1g_2,g_3)-\varphi(g_1,g_2g_3)-\varphi(g_1,g_2)
\]
（$\varphi$ 为 2-余链），以及
\footnote{译者注：原书给出的这个平凡作用特例不由式 (22.5.50) 直接化出：按该式还应
出现 $\varphi(g_2,g_3)$，中间两项的符号也相反。译文照录特例公式，不静默改写。}
\[
 \delta\varphi(g_1,g_2)=\varphi(g_1g_2)-\varphi(g_1)-\varphi(g_2)
\]
（$\varphi$ 为 1-余链）。直接计算，或利用形式表达式
$\operatorname{Def}_\mu=\delta\mu$，可知 $\operatorname{Def}_\mu$ 是 2-余闭链。

然而，当群 $G$ 非紧，例如 $G=\widetilde{\Ham}(M,\omega)$ 或类似情形时，
$(C^*(G,\mathbb R),\delta)$ 有由\emph{有界余链}组成的子复形
$(C_{\mathrm{bd}}^*(G,\mathbb R),\delta)$。当缺陷有界时，它定义
$(C_{\mathrm{bd}}^*(G,\mathbb R),\delta)$ 的一个 2-余闭链。由于 $\mu$ 本身未必有界，
$\operatorname{Def}_\mu=\delta\mu$ 并不意味着 $\operatorname{Def}_\mu$ 定义平凡的
有界 2-上同调类，尽管形式上它是余边界。

\begin{definition}
\label{def:22-homogeneous-quasimorphism}
若拟同态满足
\begin{equation}
 \mu(g^n)=n\mu(g)
 \label{eq:22-homogeneous-quasimorphism}
\end{equation}
对所有 $n\in\mathbb Z$ 及 $g\in G$ 成立，则称它\emph{齐次}。
\end{definition}

给定群 $G$ 上的拟同态 $\mu$，有一个自然过程：取渐近平均，将 $\mu$ 齐次化。需要
指出，取渐近平均在数学许多领域中都是常用做法，尤其在动力系统与遍历理论中。

\begin{definition}[齐次化]
\label{def:22-homogenization}
拟同态 $\mu:G\to\mathbb R$ 的\emph{齐次化} $\mu^{\mathrm{homo}}$ 定义为
\begin{equation}
 \mu^{\mathrm{homo}}(g)=\lim_{n\to\infty}\frac{\mu(g^n)}n,
 \label{eq:22-homogenization}
\end{equation}
其中 $g\in G$。
\end{definition}

\begin{proposition}
\label{prop:22-homogenization}
设 $\mu:G\to\mathbb R$ 为拟同态。则 $\mu^{\mathrm{homo}}$ 仍为拟同态且齐次。此外，
\begin{equation}
 |\mu^{\mathrm{homo}}(g)-\mu(g)|\leq D(\mu).
 \label{eq:22-homogenization-bound}
\end{equation}
\end{proposition}

\begin{proof}
良定性的证明本质上与引理~\ref{lem:22-asymptotic-limit} 相同，故省略。一旦良定性成立，
齐次性立即得到。拟同态的任何有界扰动仍是拟同态，所以拟同态性质将由
\eqref{eq:22-homogenization-bound} 推出。只需证明该不等式。

依定义，
\[
 \mu^{\mathrm{homo}}(g)-\mu(g)
 =\lim_{n\to\infty}\frac{\mu(g^n)}n-\mu(g)
 =\lim_{n\to\infty}\frac{\mu(g^n)-n\mu(g)}n.
\]
可改写为
\[
 \mu(g^n)-n\mu(g)
 =\sum_{j=1}^n\bigl(\mu(g^j)-\mu(g^{j-1})-\mu(g)\bigr).
\]
和式中每一项的绝对值都不超过 $D(\mu)$，故
\[
 |\mu(g^n)-n\mu(g)|\leq nD(\mu).
\]
除以 $n$ 即得~\eqref{eq:22-homogenization-bound}。
\end{proof}
\subsection{\texorpdfstring{$\widetilde{\Ham}(M,\omega)$}{Ham(M,omega) 的万有覆盖}
上的恩托夫--波尔捷罗维奇拟同态}

现说明如何从任意幂等元 $e\in QH^*(M)$（即 $e^2=e$）的谱数
$\rho(\,\cdot\,;e)$ 得到拟同态。例如，量子上同调的单位元 $e=1$ 是这个构造的自然
候选。这一步骤由恩托夫和波尔捷罗维奇 (EnP03) 发明，后来又由麦克达夫和奥斯特罗弗
(Os03) 加以改进。

考虑次数为零的子环 $QH^{(0)}(M)\subset QH(M)$。它是域
$K=\Lambda_\omega^{(0)}$ 上带单位元 $1=PD(M)$ 的交换弗罗贝尼乌斯代数。令 $e$ 为
$QH^{(0)}(M)$ 的任意非零幂等元。考虑谱映射
\[
 \rho_e:\widetilde{\Ham}(M,\omega)\longrightarrow\mathbb R,
 \qquad \rho_e(f)=\rho(f;e).
\]
三角不等式与幂等性 $e^2=e$ 蕴含
\[
 \rho_e(fg)\leq\rho_e(f)+\rho_e(g).
\]
现在要问：$(M,\omega)$ 满足何种条件时，$\rho_e$ 满足
\begin{equation}
 \rho_e(fg)\geq\rho_e(f)+\rho_e(g)-R
 \label{eq:22-quasimorphism-lower-bound}
\end{equation}
其中常数 $R$ 与 $f,g$ 无关？这将保证 $\rho_e$ 是\emph{无限维}群
$\widetilde{\Ham}(M,\omega)$ 上的拟同态。下述定理是恩托夫和波尔捷罗维奇
(EnP03) 的定理经麦克达夫和奥斯特罗弗 (Os03) 改进后的一个小变体；这里紧随
(EnP03) 的证明。

\begin{theorem}[(EnP03)]
\label{thm:22-entov-polterovich-quasimorphism}
假设 $Q=QH^{(0)}(M)$ 作为 $K$-向量空间有分解
$Q=Q_1\oplus Q'$，其中 $Q_1,Q'$ 都是 $Q$ 的子环，且 $Q_1$ 是域。以 $e$ 表示
$Q_1$ 的单位元。则存在 $R>0$，使~\eqref{eq:22-quasimorphism-lower-bound} 成立，
因而 $\rho_e$ 在 $\widetilde{\Ham}(M,\omega)$ 上定义一个拟同态。
\end{theorem}

\begin{proof}
把 $\Lambda_\omega$ 与 $QH^*(M)$ 上的赋值 $\nu$ 分别限制到
$K=\Lambda_\omega^{(0)}$、$Q=QH^{(0)}(M)$ 和 $Q_1$ 上。

\begin{lemma}
\label{lem:22-multiplicative-valuation-norm}
函数
\[
 \zeta\in K\longmapsto e^{-\nu(\zeta)}=:|\zeta|
\]
是乘性的，即对所有 $\zeta_1,\zeta_2\in K$，
$|\zeta_1\zeta_2|=|\zeta_1||\zeta_2|$。
\end{lemma}

\begin{proof}
只需证明
$\nu(\zeta_1\zeta_2)=\nu(\zeta_1)+\nu(\zeta_2)$。当 $\zeta_1=0$ 时，
$\zeta_1\zeta_2=0$，等式平凡成立，因为两端都变成 $-\infty$。
\footnote{译者注：原书同时把 $|\zeta|$ 定义为 $e^{-\nu(\zeta)}$，又称
$\nu(0)=-\infty$；按通常扩展实数运算，这会令 $|0|=+\infty$，与“范数”不相容。
原书没有在此解决这一符号约定冲突，译文照录论证并提示。}
因此设 $\zeta_1,\zeta_2\neq0$。由于 $K$ 是域，没有零因子，故
$\zeta_1\zeta_2\neq0$。依定义，
\[
 \zeta_i=a_{i,0}T^{\nu(\zeta_i)}
 +(\zeta_i-a_{i,0}T^{\nu(\zeta_i)}),
\]
其中 $0\neq a_{i,0}\in\mathbb Q$，并且
$\nu(\zeta_i-a_{i,0}T^{\nu(\zeta_i)})>\nu(\zeta_i)$，$i=1,2$。于是
\[
 \zeta_1\zeta_2=a_{1,0}a_{2,0}T^{\nu(\zeta_1)+\nu(\zeta_2)}
 +\text{“高阶项”}.
\]
由于 $a_{1,0}a_{2,0}\neq0$，得到
$\nu(\zeta_1\zeta_2)=\nu(\zeta_1)+\nu(\zeta_2)$。
\end{proof}

\begin{proposition}
\label{prop:22-field-valuation-bound}
存在 $R'$，使每个 $b\in Q_1\setminus\{0\}$ 都满足
$\nu(b)+\nu(b^{-1})\geq R'$。
\end{proposition}

\begin{proof}
域 $Q_1$ 是域 $K$ 上的有限维向量空间，因而是 $K$ 的有限域扩张。因此 $K$ 上上述
乘性范数 $|\,\cdot\,|$ 可延拓为 $Q_1$ 上的乘性范数，仍记作
$|\,\cdot\,|$。另一方面，函数
\[
 b\in Q_1\longmapsto e^{-\nu_{Q_1}(b)}
\]
在 $K$ 上的向量空间 $Q_1$ 上定义一个范数 $\lVert\,\cdot\,\rVert$。有限维
$K$-向量空间上的所有范数等价，故存在 $r>0$，使
\[
 \lVert b\rVert\leq r|b|
\]
对所有 $b\in Q_1$ 成立。于是对 $b\neq0$，
\[
 \lVert b\rVert\,\lVert b^{-1}\rVert
 \leq r^2|b|\,|b^{-1}|=r^2.
\]
而
\[
 \lVert b\rVert\,\lVert b^{-1}\rVert
 =e^{-\nu_{Q_1}(b)}e^{-\nu_{Q_1}(b^{-1})}.
\]
结合二式，得到
\begin{equation}
 \nu_{Q_1}(b)+\nu_{Q_1}(b^{-1})\geq-2\log r
 \label{eq:22-field-valuation-estimate}
\end{equation}
对所有 $0\neq b\in Q_1$ 成立。取 $R'=-2\log r$ 即可。
\end{proof}

现在证明定理~\ref{thm:22-entov-polterovich-quasimorphism}。令
$f,g\in\widetilde{\Ham}(M,\omega)$。由三角不等式，
\[
 \rho(fg;e)\geq\rho(f;e)-\rho(g^{-1};e).
\]
把引理~\ref{lem:22-spectral-duality} 应用于 $\rho(g^{-1};e)$，得到
\begin{equation}
 \rho(fg;e)\geq\rho(f;e)+\inf_{b;\,\Pi(b,e)\neq0}\rho(g;b).
 \label{eq:22-quasimorphism-duality-bound}
\end{equation}
估计右端的下确界。写 $b=b_1+b'$，其中 $b_1\in Q_1$、$b'\in Q'$。由 (20.4.43)，
\[
 0\neq\Pi(b,e)=\pi_0(\langle b\cdot e,1\rangle)
 =\pi_0(\langle b_1\cdot e,1\rangle)=\pi_0(\langle b_1,1\rangle).
\]
故 $b_1\neq0$ 且 $\nu_{Q_1}(b_1)\leq0$。由于 $Q_1$ 是域，$b_1$ 可逆，并由
\eqref{eq:22-field-valuation-estimate} 有
\[
 \nu_{Q_1}(b_1^{-1})\geq-\nu_{Q_1}(b_1)-2\log r\geq R'.
\]
另一方面，两次应用三角不等式，得到
\begin{align*}
 \rho(g;b)&\geq\rho(g;b\cdot e)-\rho(1;e)
 =\rho(g;b_1\cdot e)-\rho(1;e)\\
 &\geq\rho(g;e)-\rho(1;b_1^{-1})-\rho(1;e),
\end{align*}
其中 $1\in\widetilde{\Ham}(M,\omega)$ 是与零哈密顿量 $0$ 相伴的常哈密顿道路所对应的
恒等元。由谱不变量的归一化公理，
$\rho(1;b_1^{-1})=-\nu(b_1^{-1})$ 且 $\rho(1;e)=-\nu(e)$。所以对所有
$\Pi(b,e)\neq0$ 的 $b$，
\[
 \rho(g;b)\geq\rho(g;e)+\nu(b_1^{-1})+\nu(e)
 \geq\rho(g;e)+R'+\nu(e).
\]
从而
\[
 \inf_{b;\,\Pi(b,e)\neq0}\rho(g;b)
 \geq\rho(g;e)+R'+\nu(e).
\]
代入~\eqref{eq:22-quasimorphism-duality-bound}，得到
\[
 \rho(fg;e)\geq\rho(f;e)+\rho(g;e)+R'+\nu(e).
\]
因此 $\rho(\,\cdot\,;e)$ 定义拟同态，其缺陷小于
$R=|R'+\nu(e)|$。这就完成了定理~22.5.11 的证明。
\end{proof}

对上述拟同态 $\rho_e$ 作齐次化，得到 $\widetilde{\Ham}(M,\omega)$ 上的齐次拟同态
\begin{equation}
 \widetilde\mu_e(f)=\lim_{m\to\infty}\frac{\rho_e(f^m)}m.
 \label{eq:22-homogeneous-calabi-quasimorphism}
\end{equation}

\begin{remark}
\label{rem:22-enp-sign-convention}
需要指出，由于 (EnP03) 采用不同约定，(EnP03) 式 (17) 中的负号没有出现在这里的
定义中。
\end{remark}

现在证明拟同态 $\widetilde\mu$ 的下述重要性质。

\begin{proposition}
\label{prop:22-calabi-property}
设 $U\subset M$ 可移，即存在 $\phi\in\Ham(M,\omega)$，使
$\phi(\overline U)\cap\overline U=\varnothing$。则对所有满足
$f=[\phi_F]\in\widetilde{\Ham}(M,\omega)$、$\operatorname{supp}F\subset U$ 的 $f$，
\[
 \widetilde\mu(f)=\operatorname{Cal}_U(F).
\]
\end{proposition}

\begin{proof}
记 $\lambda=\phi_F$。由一般哈密顿道路 $\lambda$ 的 $\rho(\lambda;1)$ 定义，
\[
 \rho(\lambda^m;1)=\rho(\operatorname{Dev}(\lambda^m);1).
\]
（一般哈密顿道路的 $\operatorname{Dev}(\mu)$ 定义见定义 2.3.15。）而
\begin{align*}
 \operatorname{Dev}(\lambda^m)
 &=\underline{F^{\#m}}\\
 &=F^{\#m}-\frac1{\operatorname{vol}_\omega(M)}\int_M F_t^{\#m}\,d\mu_\omega\\
 &=F^{\#m}-\frac1{\operatorname{vol}_\omega(M)}
   \int_M\left(\sum_{k=1}^mF(t,((\phi_F^t)^{-1})^{k-1}(x))\right)d\mu_\omega\\
 &=F^{\#m}-\frac m{\operatorname{vol}_\omega(M)}\int_MF(t,x)\,d\mu_\omega,
\end{align*}
最后一个等号成立是因为对所有 $k=1,\ldots,m$，
$((\phi_F^t)^{-1})^{k-1}$ 都保持测度 $\mu_\omega$。因此
\begin{equation}
 \rho(\operatorname{Dev}(\lambda^m);1)
 =\rho(F^{\#m};1)+m\operatorname{Cal}(F).
 \label{eq:22-development-calabi}
\end{equation}
注意 $\operatorname{supp}F^{\#m}\subset U$，因为 $\operatorname{supp}F\subset U$。
选取任意满足 $\phi=\phi_H^1$ 的哈密顿量 $H$。由假设
$\phi(U)\cap\overline U=\varnothing$，命题 22.3.7 蕴含
\[
 \rho(F^{\#m};1)\leq\gamma(H).
\]
代入~\eqref{eq:22-development-calabi}，除以 $m$，再令 $m\to\infty$，得到
\[
 \widetilde\mu(f)\leq\operatorname{Cal}(F).
\]
类似地，以 $\lambda^{-1}$ 代替 $\lambda$，得到
\[
 \widetilde\mu(f^{-1})\leq\operatorname{Cal}(\overline F)
 =-\operatorname{Cal}(F).
\]
另一方面，由 $\widetilde\mu$ 的齐次性，
$\widetilde\mu(f^{-1})=-\widetilde\mu(f)$。结合二式，得到
$-\widetilde\mu(f)\leq-\operatorname{Cal}(F)$，等价于
$\widetilde\mu(f)\geq\operatorname{Cal}(F)$。故
$\widetilde\mu(f)=\operatorname{Cal}(F)$，这证明了 $e=1$ 时的命题。

同一证明也证明一般幂等元 $e$ 的第二条陈述。
\footnote{译者注：原书命题 22.5.15 只陈述了 $\widetilde\mu$ 的一个恒等式，却在证明
末称还证明了“一般幂等元 $e$ 的第二条陈述”；相应陈述似在排版中遗漏。译文保留这句
并说明缺口。}
\end{proof}

\begin{remark}
\label{rem:22-calabi-endpoint}
第 2.5.1 节已经证明，在任意正合（开）辛流形上，
$\operatorname{Cal}^{\mathrm{path}}(\lambda)$ 可表示为
\begin{equation}
 \operatorname{Cal}^{\mathrm{path}}(\lambda)
 =\frac1{n+1}\int_M(\phi^*\alpha-\alpha)\wedge\alpha\wedge\omega^{n-1}.
 \label{eq:22-calabi-endpoint-formula}
\end{equation}
特别地，$\operatorname{Cal}^{\mathrm{path}}(\lambda)$ 只依赖于终点
$\lambda(1)=\phi$，因而定义非平凡同态
\[
 \operatorname{Cal}:\Ham^c(M,d\alpha)\longrightarrow\mathbb R.
\]
把这一构造应用于可缩 $U$ 的 $\mathcal P_U^{\mathrm{ham}}$，可对每个可缩 $U$ 定义
$\operatorname{Cal}_U:\Ham^c(U,\omega)\to\mathbb R$。于是上述证明表明，\emph{只要}
$\widetilde\mu_e$ 下降为 $\Ham(M,\omega)$ 上的拟同态，就有
\[
 \mu_e(\phi)=\operatorname{Cal}_U(\phi)
\]
对所有 $\phi\in\Ham^c(U)\subset\Ham(M,\omega)$ 成立。
\end{remark}

下面给出使 $\widetilde\mu_e$ 下降到 $\Ham(M,\omega)$ 的一个简单充分条件。

\begin{proposition}[(EnP03) 的命题 3.4]
\label{prop:22-descend-calabi-quasimorphism}
若 $\pi_1(\Ham(M,\omega))$ 有限，则拟同态 $\widetilde\mu_e$ 下降为
$\Ham(M,\omega)$ 上的卡拉比拟同态。
\end{proposition}

\begin{proof}
注意 $\pi_1(\Ham(M,\omega))$ 正是覆盖投影
$\widetilde{\Ham}(M,\omega)\to\Ham(M,\omega)$ 的核。

先断言 $\pi_1(\Ham(M,\omega))\subset\widetilde{\Ham}(M,\omega)$ 位于
$\widetilde{\Ham}(M,\omega)$ 的中心。只需对任意
$g\in\pi_1(\Ham(M,\omega))$ 与 $h\in\widetilde{\Ham}(M,\omega)$ 证明
$ghg^{-1}h^{-1}=id$。把 $g,h$ 分别表示为 $g=[\lambda]$、$h=[\delta]$，其中
$\lambda,\delta:[0,1]\to\Ham(M,\omega)$ 为光滑道路，
$\lambda(1)=\lambda(0)=id$、$\delta(0)=id$。对每个 $s\in[0,1]$，
\[
 \lambda(1)\delta(s)\lambda(1)^{-1}\delta(s)^{-1}=id.
\]
故当 $s\in[0,1]$ 时，道路
\[
 t\longmapsto\lambda(t)\delta(st)\lambda(t)^{-1}\delta(st)^{-1}
\]
给出以 $id$ 为基点的环路同伦。依定义，$s=0$ 时是平凡环路，$s=1$ 时是表示
$\widetilde{\Ham}(M,\omega)$ 中 $ghg^{-1}h^{-1}$ 的道路
$t\mapsto\lambda(t)\delta(t)\lambda(t)^{-1}\delta(t)^{-1}$。
\footnote{译者注：原书在这两处道路公式中把逆元印作常值
$\lambda(1)^{-1}=id$；那样 $s=1$ 的道路并不表示群交换子。译文按所声称的
$ghg^{-1}h^{-1}$ 及标准中心性同伦改为 $\lambda(t)^{-1}$。}
故
$ghg^{-1}h^{-1}=id$，即 $gh=hg$，断言得证。

结合中心性与拟同态性质，对所有 $m$ 有
\[
 |\widetilde\mu_e(g^m)+\widetilde\mu_e(h^m)
 -\widetilde\mu_e((gh)^m)|
 =|\widetilde\mu_e(g^m)+\widetilde\mu_e(h^m)
 -\widetilde\mu_e(g^mh^m)|\leq R.
\]
另一方面，由 $\widetilde\mu_e$ 的齐次性，
$\widetilde\mu_e((gh)^m)=m\widetilde\mu_e(gh)$ 且
$\widetilde\mu_e(g^m)+\widetilde\mu_e(h^m)
=m(\widetilde\mu_e(g)+\widetilde\mu_e(h))$。代入得
\[
 m|\widetilde\mu_e(g)+\widetilde\mu_e(h)-\widetilde\mu_e(gh)|\leq R.
\]
这对所有 $m$ 成立，故
\begin{equation}
 \widetilde\mu_e(g)+\widetilde\mu_e(h)=\widetilde\mu_e(gh).
 \label{eq:22-central-additivity}
\end{equation}
反复令 $h=g$ 应用该恒等式，得到
$k\widetilde\mu_e(g)=\widetilde\mu_e(g^k)$ 对所有 $k\in\mathbb N$ 成立。但
$\pi_1(\Ham(M,\omega))$ 有限，故某个 $k_0\neq0$ 满足 $g^{k_0}=id$，于是
$\widetilde\mu_e(g^{k_0})=\widetilde\mu_e(id)=0$。这蕴含
$\widetilde\mu_e(g)=0$ 对所有 $g\in\pi_1(\Ham(M,\omega))$ 成立。

代回~\eqref{eq:22-central-additivity}，得到对所有
$g\in\pi_1(\Ham(M,\omega))$ 均有
$\widetilde\mu_e(gh)=\widetilde\mu_e(h)$。因此 $\widetilde\mu_e$ 下降为映射
$\mu_e:\Ham(M,\omega)\to\mathbb R$；定义为
$\mu_e(\phi)=\widetilde\mu_e(h)$，其中
$h\in\widetilde{\Ham}(M,\omega)$ 是 $\phi$ 的任意一个提升。

\begin{exercise}
\label{exer:22-calabi-quasimorphism-descent}
证明上述定义的 $\mu_e:\Ham(M,\omega)\to\mathbb R$ 具有卡拉比拟同态性质，
从而完成证明。
\end{exercise}
\end{proof}

为了稍后在第 22.6 节把拟同态延拓到拓扑哈密顿情形，现在把拟同态
$\widetilde\mu_e$ 提升到道路空间
$\mathcal P^{\mathrm{ham}}(\Symp(M,\omega),id)$，而不是万有覆盖空间
$\widetilde{\Ham}(M,\omega)$ 上。在光滑情形中，由于 $\rho(\,\cdot\,;e)$ 的同伦不变性，
这一提升在下述定义下自动成立。

\begin{definition}
\label{def:22-path-homogeneous-quasimorphism}
用同一公式定义齐次拟同态
\[
 \mu^{\mathrm{path}}=\mu_1^{\mathrm{path}}:
 \mathcal P^{\mathrm{ham}}(\Symp(M,\omega),id)\longrightarrow\mathbb R,
\]
即\footnote{译者注：原书式 (22.5.61) 的极限下标印作 $i\to\infty$，但被积式与分母
均使用 $m$；译文据此改为 $m\to\infty$。}
\begin{equation}
 \mu^{\mathrm{path}}(\lambda)
 =\lim_{m\to\infty}\frac{\rho(\lambda^m;1)}m.
 \label{eq:22-path-homogeneous-quasimorphism}
\end{equation}
\end{definition}

对任意给定开子集 $U\subset M$，仍以
\[
 \mathcal P^{\mathrm{ham}}(\Symp_U(M,\omega),id)
\]
表示支集在 $U$ 中的哈密顿道路集合。回忆
$\operatorname{Cal}_U(F)$ 由积分
\begin{equation}
 \operatorname{Cal}_U(F)
 =\frac1{\operatorname{vol}_\omega(M)}\int_0^1\int_MF(t,x)\,d\mu_\omega\,dt
 \label{eq:22-calabi-integral}
\end{equation}
定义。
\footnote{译者注：原书式 (22.5.62) 在 $F(t,x)$ 后只印 $\mu_\omega$ 而没有微分号；
译文按本书此前卡拉比积分的写法补作 $d\mu_\omega$。}
把 $\widetilde{\Ham}_U(M,\omega)$ 换成
$\mathcal P^{\mathrm{ham}}(\Symp_U(M,\omega),id)$，逐字重复上述构造，则
$\widetilde\mu_e$ 的全部上述性质对 $\mu_e^{\mathrm{path}}$ 同样成立。

本节最后讨论二维球面 $S^2$ 的谱直径。

\begin{example}
\label{ex:22-sphere-spectral-diameter}
考虑谱函数
\[
 \rho_1:\widetilde{\Ham}(M,\omega)\longrightarrow\mathbb R,
 \qquad \rho_1(f)=\rho(f;1).
\]
显然它满足 $\rho_1(fg)\leq\rho_1(f)+\rho_1(g)$。注意 $S^2$ 与量子上同调类
$1\in QH^0(S^2)$ 满足定理~\ref{thm:22-entov-polterovich-quasimorphism} 所需的假设，
因此存在与 $f,g$ 无关的常数 $R$，使
\begin{equation}
 \rho_1(fg)\geq\rho_1(f)+\rho_1(g)-R.
 \label{eq:22-sphere-quasimorphism-bound}
\end{equation}
现在显式计算这个界 $R$。

先注意，式~\eqref{eq:22-spectral-duality} 的右端变为
\[
 -\inf_{a\in QH^0(S^2)\setminus\{0\}}\rho(g;a).
\]
此外，$QH^0(S^2)$ 分解为
\[
 QH^0(S^2)=\mathbb Q1\oplus\mathbb Q\{[\omega]q^{-1}\},
\]
其中 $1=PD[S^2]$，$q$ 是次数为 $2$、赋值 $\nu(q)=4\pi$ 的形式参数，而且
$QH^0(S^2)$ 的每个非零元素都可逆。

对任意可逆的 $a\in QH^0(S^2)\setminus\{0\}$，由谱不变量归一化公理
$\rho(1;a^{-1})=-\nu(a^{-1})$，
\[
 \rho(g;a)\geq\rho(g;a\cdot a^{-1})-\rho(1;a^{-1})
 =\rho(g;1)+\nu(a^{-1}).
\]
而
\begin{equation}
 \inf_{b\in QH^0(S^2)\setminus\{0\}}\nu(b^{-1})=-4\pi.
 \label{eq:22-sphere-valuation-bound}
\end{equation}
故对所有 $a\in QH^0(S^2)\setminus\{0\}$，
$\rho(g;a)\geq\rho(g;1)-4\pi$。

令 $f,g\in\widetilde{\Ham}(M,\omega)$。由三角不等式，
\[
 \rho(fg;1)\geq\rho(f;1)-\rho(g^{-1};1).
\]
应用引理~\ref{lem:22-spectral-duality} 与~\eqref{eq:22-sphere-valuation-bound}，得到
\[
 \rho(fg;1)\geq\rho(f;1)+\rho(g;1)-4\pi.
\]
取 $f=g^{-1}$，便以 $R=4\pi$ 完成证明。综上得到下述命题。
\end{example}

\begin{proposition}
\label{prop:22-sphere-inverse-spectral-bound}
对面积为 $4\pi$ 的 $S^2$ 上任意哈密顿量 $H$，
\begin{equation}
 \rho((\phi_H^1)^{-1};1)\leq-\rho(\phi_H;1)+4\pi.
 \label{eq:22-sphere-inverse-bound}
\end{equation}
\end{proposition}

上述讨论的一个有趣推论给出 $\Ham(S^2)$ 谱直径的如下界。

\begin{theorem}
\label{thm:22-sphere-spectral-diameter}
以 $\operatorname{Diam}_\rho$ 表示\emph{谱直径}，定义为
\[
 \operatorname{Diam}_\rho(\Ham(S^2))
 :=\sup_H\{\gamma(\phi_H)\mid H\in C^\infty([0,1]\times S^2,\mathbb R)\}.
\]
则 $\operatorname{Diam}_\rho(\Ham(S^2))\leq4\pi$；等价地，
\begin{equation}
 \gamma(\phi_H)\leq4\pi\;(=\operatorname{Area}(S^2))
 \label{eq:22-sphere-spectral-diameter}
\end{equation}
对所有 $H$ 成立。
\footnote{译者注：原书在 $\operatorname{Diam}_\rho(\Ham(S^2))$ 的定义中把哈密顿量
的定义域印作 $[0,1]\times M$，但本定理从头至尾讨论 $S^2$；译文据此改作
$[0,1]\times S^2$。}
\end{theorem}

$S^2$ 谱直径的这个上界与波尔捷罗维奇关于 $S^2$ 的霍费尔直径的定理
（(Po98b) 中的定理 17.6.3）形成鲜明对比：后者断言相对于霍费尔距离，
$\Ham(S^2)$ 的直径为无穷。事实上，只要量子上同调 $QH^*(M)$ 成为域，例如
$\mathbb{CP}^n$，类似界就成立。

由此自然产生一个有趣问题。

\begin{question}
\label{ques:22-hofer-spectral-diameter}
找出 $(M,\omega)$ 具有有限（或无限）霍费尔直径的精确判据。谱直径又如何？这对
哈密顿动力学有何含义？
\end{question}

\section{回到拓扑哈密顿动力学}

第 6 章引入了拓扑哈密顿流及其相伴哈密顿量的概念。本节说明如何把各种谱量延拓到
\emph{拓扑}哈密顿道路。

\subsection{拓扑哈密顿量的谱不变量}

本小节把\emph{哈密顿道路}的谱不变量定义及基本性质延拓到拓扑哈密顿范畴。

为完成这一延拓，关键是像 (Oh05c) 那样在哈密顿道路空间
\[
 \mathcal P^{\mathrm{ham}}(\Symp(M,\omega),id)
\]
的层面，而不只在
$\Ham(M,\omega)$ 的覆盖空间上给出定义。此外，第 6.3 节的拓扑哈密顿量唯一性定理
对延拓到这一 $C^0$ 范畴至关重要。

先用哈密顿道路 $\phi_H$ 改写谱数 $\rho(H;a)$：
\[
 \rho(\lambda;a):=\rho(\underline H;a)
 =\inf_{\alpha\in\ker\partial_H}\{\lambda_{\underline H}(\alpha)\mid[\alpha]=a^\flat\}.
\]
对光滑哈密顿量 $H$，这只是换一种写法。为了用一列 $H_i$（其哈密顿流
$\phi_{H_i}$ 一致收敛到给定的拓扑哈密顿道路 $\lambda$）所对应的
$\rho(\underline H_i;a)$ 之极限来延拓定义，需要证明
$\lim_{i\to\infty}\rho(\underline H_i;a)$ 与这类序列 $H_i$ 的选择无关。

\begin{theorem}
\label{thm:22-topological-spectral-invariant}
对带归一化哈密顿量 $H$ 的光滑哈密顿道路 $\phi_H$，定义
\[
 \rho(\phi_H;a)=\rho(\underline H;a).
\]
则映射 $\rho_a:\phi_H\mapsto\rho(\phi_H;a)$ 延拓为（在哈密顿拓扑下）连续函数
\[
 \rho_a=\rho(\,\cdot\,;a):
 \mathcal P^{\mathrm{ham}}(\operatorname{Sympeo}(M,\omega),id)
 \longrightarrow\mathbb R,
\]
并满足三角不等式
\begin{equation}
 \rho(\lambda\mu;a\cdot b)\leq\rho(\lambda;a)+\rho(\mu;b).
 \label{eq:22-topological-triangle}
\end{equation}
\end{theorem}

\begin{proof}
第一条陈述是
\[
 \rho_a:\mathcal P^{\mathrm{ham}}(\Symp(M,\omega),id)\longrightarrow\mathbb R
\]
的哈密顿连续性与定理 6.3.3 的直接推论：对任意拓扑哈密顿道路 $\lambda$，定义
\begin{equation}
 \rho(\lambda;a)=\lim_{i\to\infty}\rho(\underline H_i;a),
 \label{eq:22-topological-spectral-limit}
\end{equation}
其中 $H_i$ 是任意满足 $\phi_{H_i}\to\lambda$ 且在 $L^{(1,\infty)}$ 中收敛的序列。
唯一性定理保证该定义良定。此外，命题 21.3.6 蕴含此延拓在
$\mathcal P^{\mathrm{ham}}(\operatorname{Sympeo}(M,\omega),id)$ 上连续。对光滑情形的
三角不等式取极限，立刻得到~\eqref{eq:22-topological-triangle}。
\end{proof}

现在专注于 $1\in QH^*(M)$ 的不变量 $\rho(\lambda;1)$。回忆，对光滑哈密顿道路
$\lambda=\phi_H$，此前定义了
\[
 \gamma(H)=\rho(H;1)+\rho(\overline H;1).
\]
为免混淆，当把 $\gamma$ 看作
$\mathcal P^{\mathrm{ham}}(\Symp(M,\omega),id)$ 上的函数时，记作
$\operatorname{norm}_\gamma$。

\begin{definition}[谱伪范数]
\label{def:22-spectral-pseudonorm}
令 $\lambda\in\mathcal P^{\mathrm{ham}}(\Symp(M,\omega),id)$，并令 $H$ 为满足
$\lambda=\phi_H$ 的哈密顿量。定义函数
\[
 \operatorname{norm}_\gamma:
 \mathcal P^{\mathrm{ham}}(\Symp(M,\omega),id)\longrightarrow\mathbb R_+
\]
为 $\operatorname{norm}_\gamma(\lambda)=\gamma(H)$。称
$\operatorname{norm}_\gamma(\lambda)$ 为 $\lambda$ 的\emph{谱伪范数}。
\end{definition}

此前已经证明它非负，并且只依赖于 $\lambda=\phi_H$ 在 $\Ham(M,\omega)$ 中的道路同伦类。
拓扑哈密顿量的唯一性再次使我们能把定义延拓到拓扑哈密顿道路。

\begin{proposition}
\label{prop:22-topological-spectral-pseudonorm}
谱伪范数 $\gamma$ 延拓为连续函数
\[
 \operatorname{norm}_{\overline\gamma}:
 \mathcal P^{\mathrm{ham}}(\operatorname{Sympeo}(M,\omega),id)
 \longrightarrow\mathbb R,
\]
其定义为
\begin{equation}
 \operatorname{norm}_{\overline\gamma}(\lambda)
 =\lim_{i\to\infty}\operatorname{norm}_\gamma(\phi_{H_i}),
 \label{eq:22-topological-spectral-pseudonorm}
\end{equation}
其中 $\phi_{H_i}\to\lambda$ 是 $L^{(1,\infty)}$-哈密顿拓扑中的某个（因而任意）序列。
\end{proposition}

回忆，对光滑哈密顿量 $H$，每个
$\rho(\phi_H;a)=\rho(H;a)$ 都与哈密顿方程
$\dot x=X_H(t,x)$ 的一个周期轨道相伴，并对应于该周期轨道的作用量；至少这对有理辛
流形上的\emph{任意}哈密顿量成立（证明见 (Oh05c, Oh09a)），也对一般无理辛流形上的\emph{非退化}
哈密顿量成立（见 (Oh09a), (Ush08)）。由此看来，下述问题似乎至关重要。

\begin{question}
\label{ques:22-topological-pseudonorm-dynamics}
\stfhIndex{I178}{Puweifanshu}{谱伪范数（spectral pseudonorm）}
延拓后的谱伪范数 $\operatorname{norm}_{\overline\gamma}(\lambda)$ 对拓扑哈密顿流的动力学
有何含义？
\end{question}

下面是 (Oh05d) 中谱位移能定义在拓扑哈密顿范畴内的对应版本。

\begin{definition}[谱位移能]
\label{def:22-topological-spectral-displacement-energy}
令 $A\subset M$ 为紧子集。定义 $A$ 的谱位移能 $e_{\overline\gamma}(A)$ 为
\[
 e_{\overline\gamma}(A)
 =\inf_\lambda\{\operatorname{norm}_{\overline\gamma}(\lambda)
 \mid A\cap\lambda(1)(A)=\varnothing,
 \ \lambda\in\mathcal P_\infty^{\mathrm{ham}}
 (\operatorname{Sympeo}(M,\omega),id)\}.
\]
\end{definition}

展开哈密顿同胚与谱位移能的定义，还得到下述定理。

\begin{theorem}
\label{thm:22-spectral-displacement-agreement}
对任意紧子集 $A\subset M$，有
$e_{\overline\gamma}(A)=e_\gamma(A)$。
\end{theorem}

基于该定理，即使在拓扑哈密顿范畴中，也把 $A$ 的谱位移能简记为 $e_\gamma(A)$。于是
有下述定理。

\begin{theorem}
\label{thm:22-spectral-displacement-invariance}
对每个 $\psi\in\operatorname{Sympeo}(M,\omega)$，
\[
 e_\gamma(A)=e_\gamma(\psi(A)).
\]
\end{theorem}

\begin{proof}
注意 $h(A)\cap A=\varnothing$ 当且仅当
$\psi h\psi^{-1}(\psi(A))\cap\psi(A)=\varnothing$。此外，
$h\in\Hameo(M,\omega)$ 当且仅当
$\psi h\psi^{-1}\in\Hameo(M,\omega)$。
\end{proof}

接着考虑谱范数
\[
 \gamma(\phi):=\inf_{\operatorname{ev}_1(\lambda)=\phi}
 \operatorname{norm}_\gamma(\lambda)
 =\inf_{H\mapsto\phi}\{\rho(H;1)+\rho(\overline H;1)\}.
\]
下述定义把谱范数延拓到 $\Hameo(M,\omega)$。

\begin{definition}[谱范数]
\label{def:22-generalized-spectral-norm}
令 $h\in\Hameo(M,\omega)$，考虑满足
$\overline{\operatorname{ev}}_1(\lambda)=h$ 的拓扑哈密顿道路
$\lambda\in\mathcal P_\infty^{\mathrm{ham}}
(\operatorname{Sympeo}(M,\omega),id)$。若
$\overline{\operatorname{ev}}_1(\lambda)=h$，简记为 $\lambda\mapsto h$。定义
\footnote{译者注：原书定义首句把端点评价条件印作
$\overline{\operatorname{ev}}_1(0)=h$，但下一句及式 (22.6.70) 都使用
$\overline{\operatorname{ev}}_1(\lambda)=h$；译文按后两处改正。}
\begin{equation}
 \overline\gamma(h)=\inf_\lambda
 \{\operatorname{norm}_{\overline\gamma}(\lambda)
 \mid \lambda\in\mathcal P_\infty^{\mathrm{ham}}
 (\operatorname{Sympeo}(M,\omega),id),
 \ \overline{\operatorname{ev}}_1(\lambda)=h\}.
 \label{eq:22-generalized-spectral-norm}
\end{equation}
\end{definition}

下述结果在拓扑哈密顿语境中建立了不变范数全部性质的对应版本。

\begin{theorem}
\label{thm:22-generalized-spectral-norm-properties}
广义谱函数
$\overline\gamma:\Hameo(M,\omega)\to\mathbb R_+$ 满足不变范数的全部性质。
\end{theorem}

一旦证明光滑情形的下述连续性引理，证明便与霍费尔范数的证明本质相同。

\begin{lemma}
\label{lem:22-spectral-norm-hofer-continuity}
函数 $\gamma:\Ham(M,\omega)\to\mathbb R_+$ 在霍费尔拓扑中连续，因而也在
$\Ham(M,\omega)$ 的哈密顿拓扑中连续。
\end{lemma}

\begin{proof}
设 $H\mapsto\phi$、$K\mapsto\psi$。$\gamma$ 的三角不等式及
$\gamma(\phi)\leq\lVert\phi\rVert$ 蕴含
\[
 |\gamma(\phi)-\gamma(\psi)|
 \leq\gamma(\phi^{-1}\psi)
 \leq\lVert\phi^{-1}\psi\rVert
 =d_{\mathrm{Hofer}}(\phi,\psi).
\]
\end{proof}

由此产生如下基础问题。

\begin{question}
\label{ques:22-spectral-norm-continuity}
\begin{enumerate}[label=(\arabic*)]
\item $\gamma$（或 $\overline\gamma$）在 $C^0$-拓扑中连续吗？
\item 恒等式
\begin{equation}
 \overline\gamma|_{\Ham(M,\omega)}=\gamma
 \label{eq:22-spectral-norm-restriction}
\end{equation}
成立吗？
\end{enumerate}
\end{question}

最近，塞法迪尼 (Sey13a) 综合 $C^0$ 辛拓扑与哈密顿动力学中的多种要素，对二维曲面
肯定回答了第一个问题。它又在二维保面积动力系统中导出一个有趣推论 (Sey13b)。

\subsection{\texorpdfstring{$\mathcal P^{\mathrm{ham}}(\operatorname{Sympeo}(S^2),id)$}
{P-ham(Sympeo(S2),id)} 上的卡拉比拟同态}

本节余下部分把讨论限制在带标准辛形式 $\omega_{S^2}$ 的球面 $S^2$。在记号中省略
$\omega_{S^2}$，以
\[
 \mathcal P^{\mathrm{ham}}(\operatorname{Sympeo}(S^2),id)
 \quad\text{和}\quad \Hameo(S^2)
\]
分别简记 $S^2$ 上的拓扑哈密顿道路群与哈密顿同胚群，其他记号类似。

由上述定义及 $\rho(\,\cdot\,;1)$ 的哈密顿连续性，
$\mu^{\mathrm{path}}$ 也哈密顿连续。$\mu^{\mathrm{path}}$ 的卡拉比性质立即蕴含：
把 $\mu$ 限制到
$\mathcal P^{\mathrm{ham}}(\Symp_U(S^2),id)$ 后具有如下同态性质。

\begin{corollary}
\label{cor:22-local-path-calabi-homomorphism}
设 $U$ 是 $S^2$ 的开子集，$\overline U$ 在 $S^2$ 上可移，并令
\[
 \lambda_1,\lambda_2\in\mathcal P^{\mathrm{ham}}(\Symp_U(S^2),id).
\]
则
\[
 \mu^{\mathrm{path}}(\lambda_1\lambda_2)
 =\mu^{\mathrm{path}}(\lambda_1)+\mu^{\mathrm{path}}(\lambda_2).
\]
\end{corollary}

上述全部讨论都可延拓到拓扑哈密顿道路层面。只要知道下列事实，推广便立即得到：
\begin{enumerate}[label=(\arabic*)]
\item 对任意闭辛流形，$\rho(\,\cdot\,;1)$ 已延拓到
$\mathcal P^{\mathrm{ham}}(\operatorname{Sympeo}(M,\omega),id)$；特别地，第 12.6 节
处理了 $(S^2,\omega_{S^2})$\footnote{译者注：原书此处交叉引用为第 12.6 节；但所指的正是本节前述拓扑延拓，
且本书第 12.6 节并不讨论该问题，疑应为第 22.6 节。译文保留原编号并提示。}；
\item 此外，该延拓哈密顿连续，即在哈密顿拓扑中连续。
\end{enumerate}
综上得到下述定理。

\begin{theorem}
\label{thm:22-topological-calabi-quasimorphism}
卡拉比拟同态
\[
 \mu^{\mathrm{path}}:
 \mathcal P^{\mathrm{ham}}(\Symp(S^2),id)\longrightarrow\mathbb R
\]
可延拓为拟同态
\[
 \overline\mu^{\mathrm{path}}:
 \mathcal P^{\mathrm{ham}}(\operatorname{Sympeo}(S^2),id)
 \longrightarrow\mathbb R,
\]
并满足命题~\ref{prop:22-descend-calabi-quasimorphism} 与卡拉比性质的所有对应版本。
\end{theorem}

下面提出一个猜想；我们相信，它在研究 $S^2$ 保面积群的单纯性问题时将起重要作用。
回忆 (EnP03) 已由恩托夫和波尔捷罗维奇对 $S^2$ 上\emph{光滑}哈密顿微分同胚群
$\Ham(S^2)$ 证明了相应事实。

\begin{conjecture}
\label{conj:22-descend-topological-calabi}
令
$\overline\mu^{\mathrm{path}}:
\mathcal P^{\mathrm{ham}}(\operatorname{Sympeo}(S^2),id)\to\mathbb R$
为~\eqref{eq:22-path-homogeneous-quasimorphism} 给出的齐次卡拉比拟同态之上述延拓。
它下降为齐次拟同态
$\overline\mu:\Hameo(S^2)\to\mathbb R$，满足
\begin{equation}
 \overline\mu^{\mathrm{path}}
 =\overline\mu\circ\overline{\operatorname{ev}}_1.
 \label{eq:22-topological-calabi-descent}
\end{equation}
特别地，只要
$\lambda\in\mathcal P^{\mathrm{ham}}(\operatorname{Sympeo}(S^2),id)$，
$\overline\mu^{\mathrm{path}}(\lambda)$ 便只依赖于 $\lambda$ 的时一映射 $\lambda(1)$。
此外，对任意 $\phi\in\Hameo(S^2)$，只要存在支集在 $\operatorname{Int}(D^+)$ 中、
且满足 $\lambda(1)=\phi$ 的拓扑哈密顿道路 $\lambda$，就有
\begin{equation}
 \overline\mu(\phi)=\overline{\operatorname{Cal}}_{D^+}(\phi).
 \label{eq:22-topological-calabi-property}
\end{equation}
\end{conjecture}

猜想~\ref{conj:22-descend-topological-calabi} 与
$\overline\mu^{\mathrm{path}}$ 的卡拉比性质会立即推出下述猜想。

\begin{conjecture}
\label{conj:22-calabi-homeomorphism-extension}
卡拉比同态
$\operatorname{Cal}:\Ham(D^2,\partial D^2)\to\mathbb R$ 可延拓为在哈密顿拓扑中连续的
同态
\[
 \overline{\operatorname{Cal}}:
 \Hameo(D^2,\partial D^2)\longrightarrow\mathbb R.
\]
\end{conjecture}

\stfhIndex{I199}{Yebaomianjitongpei}{野保面积同胚（wild area-preserving homeomorphism）}
\section{\texorpdfstring{$D^2$}{D2} 上的野保面积同胚}

本节描述 $\operatorname{Sympeo}(D^2,\partial D^2)$ 中一个紧支保面积同胚的例子；若
猜想~\ref{conj:22-calabi-homeomorphism-extension} 成立，它将不属于
$\Hameo(D^2,\partial D^2)$。由定理 6.4.1，这将蕴含
$\Hameo(D^2,\partial D^2)$ 是
$\operatorname{Sympeo}(D^2,\partial D^2)$ 的真正规子群。再结合 (Oh06b), (Sik07)
中的下述光滑化定理，上述一系列陈述将导出
$\operatorname{Homeo}_\Omega(D^2,\partial D^2)$ 非单纯。

\begin{theorem}
\label{thm:22-area-preserving-smoothing}
对 $D^2$ 上的标准面积形式 $\Omega$，把它也看作辛形式 $\omega=\Omega$，有
\[
 \operatorname{Sympeo}(D^2,\partial D^2)
 =\operatorname{Homeo}_\Omega(D^2,\partial D^2).
\]
\end{theorem}

下面集中构造 $D^2$ 上野保面积同胚的例子。为此，需要考虑 $D^2$ 的重标度映射
\[
 R_a:D^2(1)\longrightarrow D^2(a)\subset D^2(1),
 \qquad 0<a<1,
\]
在 $\Hameo(D^2,\partial D^2)$ 上的共轭作用；这里 $D^2(a)$ 是以原点为中心、半径为
$a$ 的圆盘。$R_a$ 是共形辛映射，所以只要其共轭作用有定义，就把辛映射映成辛映射。

此外，右复合 $R_a$ 定义映射
\[
 \phi\longmapsto\phi\circ R_a:\Hameo(D^2,\partial D^2)
 \longrightarrow\Hameo(D^2(a),\partial D^2(a))
 \subset\Hameo(D^2,\partial D^2),
\]
而复合 $R_a^{-1}$ 定义映射
\[
 \Hameo(D^2(a),\partial D^2(a))
 \longrightarrow\Hameo(D^2,\partial D^2).
\]
关于亚历山大同痕下卡拉比不变量的行为，有如下重要公式。

\begin{lemma}
\label{lem:22-calabi-rescaling}
令 $\lambda$ 为 $D^2$ 上给定的拓扑哈密顿道路，并定义映射 $\lambda_a$ 为
\[
 \lambda_a(t,x)=
 \begin{cases}
  a\lambda(t,x/a),&|x|\leq a(1-\eta),\\
  x,&\text{其他情形},
 \end{cases}
 \qquad 0<a\leq1.
\]
则 $\lambda_a$ 也是 $D^2$ 上的拓扑哈密顿道路，并满足
\begin{equation}
 \overline{\operatorname{Cal}}^{\mathrm{path}}(\lambda_a)
 =a^4\overline{\operatorname{Cal}}^{\mathrm{path}}(\lambda).
 \label{eq:22-calabi-rescaling}
\end{equation}
\end{lemma}

\begin{proof}
直接计算可知，$\lambda_a$ 由如下（唯一）拓扑哈密顿量生成：
\[
 \operatorname{Dev}(\lambda_a)(t,x)=
 \begin{cases}
  a^2H(t,x/a),&|x|\leq a(1-\eta),\\
  0,&\text{其他情形},
 \end{cases}
\]
其中 $H=\operatorname{Dev}(\lambda)$。显然，若 $\lambda_i$ 是 $\lambda$ 的一列光滑
哈密顿逼近，且 $\lambda_{i,a}$ 由与 $\lambda_i$ 相同的公式定义，则右端函数是
$\operatorname{Dev}(\lambda_{i,a})$ 的哈密顿极限。由此推出
\begin{align*}
 \overline{\operatorname{Cal}}^{\mathrm{path}}(\lambda_a)
 &=\int_0^1\int_{D^2(a(1-\eta))}a^2H\!\left(t,\frac xa\right)
   \Omega\wedge dt\\
 &=a^4\int_0^1\int_{D^2}H(t,y)\Omega\wedge dt
 =a^4\overline{\operatorname{Cal}}^{\mathrm{path}}(\lambda).
\end{align*}
这证明了~\eqref{eq:22-calabi-rescaling}。
\end{proof}

下面构造一个野保面积同胚的例子 (Oh10)。

\begin{example}
\label{ex:22-wild-area-preserving-homeomorphism}
作好上述准备，考虑二进数集合 $1/2^k$（$k=0,\ldots$）。令 $(r,\theta)$ 为 $D^2$
上的极坐标，则标准面积形式为
\[
 \omega=r\,dr\wedge d\theta.
\]
考虑如下形式的映射 $\phi_k:D^2\to D^2$：
\[
 \phi_k=\phi_{\rho_k}:(r,\theta)\longmapsto(r,\theta+\rho_k(r)),
\]
其中 $\rho_k:(0,1]\to[0,\infty)$ 是支集在 $(0,1)$ 中的光滑函数。于是
$\phi_{\rho_k}$ 是由自治哈密顿量
\[
 F_{\phi_k}(r,\theta)=-\int_1^r s\rho_k(s)\,ds
\]
生成的保面积映射。因此其卡拉比不变量为
\begin{equation}
 \operatorname{Cal}(\phi_k)
 =-\int_{D^2}\left(\int_1^r s\rho_k(s)\,ds\right)r\,dr\,d\theta
 =\pi\int_0^1r^3\rho_k(r)\,dr.
 \label{eq:22-twist-calabi}
\end{equation}
\footnote{译者注：原书式 (22.7.75) 的末积分变量印作 $dt$，但被积函数及积分上下限
均以 $r$ 为变量；译文改作 $dr$。}
如下选取 $\rho_k$：
\begin{enumerate}[label=(\arabic*)]
\item $\rho_k$ 的支集在 $1/2^k<r<1/2^{k-1}$ 中；
\item 对每个 $k=1,\ldots$，当 $r\in(1/2^k,1/2^{k-1})$ 时，
\begin{equation}
 \rho_k(r)=2^4\rho_{k-1}(2r);
 \label{eq:22-twist-rescaling}
\end{equation}
\item $\operatorname{Cal}(\phi_1)=1$。
\end{enumerate}

按构造，$\phi_k$ 的支集彼此不交，故可以自由复合而无需考虑复合次序。因此无穷乘积
\[
 \prod_{k=0}^\infty\phi_k
\]
\footnote{译者注：原书这里把乘积下标写成 $k=0$，但紧接着的分段定义、后文
$\psi_N$ 与定理证明均从 $k=1$ 开始；并且按性质 (1)，$\rho_0$ 在其定义域
$(0,1]$ 内不可能支撑于 $1<r<2$。这里疑应为 $\prod_{k=1}^\infty\phi_k$；译文保留
原下标并提示索引不一致。}
良定，并定义一个除原点外光滑的连续映射；在原点处，$\phi_\rho$ 连续但不可微。
这个无穷乘积也可写成同胚 $\phi_\rho$，其取值为 $\phi_\rho(0)=0$ 以及
\[
 \phi_\rho(r,\theta)=(r,\theta+\rho(r)),
\]
其中光滑函数 $\rho:(0,1]\to\mathbb R$ 定义为
\[
 \rho(r)=\rho_k(r),\qquad
 r\in[1/2^k,1/2^{k-1}],\quad k=1,2,\ldots.
\]
容易验证 $\phi_\rho$ 在 $D^2\setminus\{0\}$ 上光滑，在 $0$ 处也连续，并与上述无穷
乘积一致。显然 $\phi_{-\rho}$ 是 $\phi_\rho$ 的逆，故它是同胚。此外，
\[
 \phi_\rho^*(r\,dr\wedge d\theta)=r\,dr\wedge d\theta
 \qquad\text{于 }D^2\setminus\{0\},
\]
所以 $\phi_\rho$ 确实保面积。
\end{example}

下述引理将在定理~\ref{thm:22-wild-not-hameomorphism} 的证明中起重要作用。

\begin{lemma}
\label{lem:22-twist-self-similarity}
令 $\phi_k$ 为例~\ref{ex:22-wild-area-preserving-homeomorphism} 中的微分同胚。则
\begin{equation}
 R_{\frac12}\circ\phi_{k-1}^{2^4}\circ R_{\frac12}^{-1}=\phi_k.
 \label{eq:22-twist-conjugacy}
\end{equation}
特别地，
\begin{equation}
 \operatorname{Cal}(\phi_k)=\operatorname{Cal}(\phi_{k-1}).
 \label{eq:22-twist-calabi-constant}
\end{equation}
\end{lemma}

\begin{proof}
由~\eqref{eq:22-twist-rescaling}，计算得
\[
 R_{\frac12}\circ\phi_{k-1}\circ R_{\frac12}^{-1}(r,\theta)
 =(r,\theta+\rho_{k-1}(2r))
 =\left(r,\theta+\frac1{2^4}\rho_k(r)\right).
\]
迭代该恒等式 $2^4$ 次，便由~\eqref{eq:22-twist-rescaling} 得到
\eqref{eq:22-twist-conjugacy}。再由它与~\eqref{eq:22-calabi-rescaling} 得到
\eqref{eq:22-twist-calabi-constant}。
\end{proof}

该引理与式~\eqref{eq:22-twist-rescaling} 立即给出如下推论。
\footnote{译者注：推论实际还使用例~\ref{ex:22-wild-area-preserving-homeomorphism}
中的初值 $\operatorname{Cal}(\phi_1)=1$，且递推恒等式直接来自
式~\eqref{eq:22-twist-calabi-constant}；原书这里只引用式 (22.7.76)。}

\begin{corollary}
\label{cor:22-twist-calabi-one}
对所有 $k=1,\ldots$，
\[
 \operatorname{Cal}(\phi_k)=1.
\]
\end{corollary}

现在可以证明下述定理。

\begin{theorem}
\label{thm:22-wild-not-hameomorphism}
若猜想~\ref{conj:22-calabi-homeomorphism-extension} 成立，则
$\phi_\rho$ 不属于 $\Hameo(D^2,\partial D^2)$。
\end{theorem}

\begin{proof}
反设 $\phi_\rho\in\Hameo(D^2,\partial D^2)$，即存在道路
$\lambda\in\mathcal P^{\mathrm{ham}}
(\operatorname{Sympeo}(D^2,\partial D^2),id)$，满足 $\lambda(1)=\phi_\rho$。
于是其卡拉比不变量有有限值，记作
\begin{equation}
 \overline{\operatorname{Cal}}(\phi_\rho)
 =\overline{\operatorname{Cal}}^{\mathrm{path}}(\lambda)=C_1
 \label{eq:22-wild-calabi-value}
\end{equation}
其中 $C_1\in\mathbb R$。

写 $\phi_\rho=\psi_N\widetilde\psi_N$，其中
\[
 \psi_N=\prod_{i=1}^N\phi_i,
 \qquad
 \widetilde\psi_N=\prod_{i=N+1}^\infty\phi_i.
\]
取 $N=1$，由 $\overline{\operatorname{Cal}}$ 的同态性质得到
\begin{equation}
 C_1=\overline{\operatorname{Cal}}(\psi_1)
 +\overline{\operatorname{Cal}}(\widetilde\psi_1).
 \label{eq:22-wild-calabi-split}
\end{equation}
注意 $\psi_N$ 光滑，故显然属于 $\Hameo(D^2,\partial D^2)$。若 $\phi_\rho$ 属于该群，
由 $\Hameo(D^2,\partial D^2)$ 的群性质，$\widetilde\psi_N$ 也属于该群。又由
推论~\ref{cor:22-twist-calabi-one}，
\[
 \overline{\operatorname{Cal}}(\psi_1)
 =\operatorname{Cal}(\phi_1)=1,
\]
所以
\begin{equation}
 \overline{\operatorname{Cal}}(\widetilde\psi_1)=C_1-1.
 \label{eq:22-wild-tail-calabi}
\end{equation}

另一方面，把~\eqref{eq:22-twist-conjugacy} 迭代应用于无穷乘积
\[
 \widetilde\psi_1=\prod_{i=2}^\infty\phi_i,
\]
可知 $\widetilde\psi_1$ 满足
\begin{equation}
 \widetilde\psi_1(r,\theta)=
 \begin{cases}
  R_{\frac12}\circ\phi_\rho^{2^4}\circ R_{\frac12}^{-1}(r,\theta),
       &0<r\leq\frac12,\\
  (r,\theta),&\frac12\leq r\leq1.
 \end{cases}
 \label{eq:22-wild-tail-rescaling}
\end{equation}
特别地，
\begin{equation}
 \widetilde\psi_1
 =\overline{\operatorname{ev}}_1((\lambda^{2^4})_{1/2}).
 \label{eq:22-wild-tail-evaluation}
\end{equation}
将引理~\ref{lem:22-calabi-rescaling} 用于 $a=1/2$，再用
$\overline{\operatorname{Cal}}^{\mathrm{path}}$ 的同态性质，得到
\begin{align}
 \overline{\operatorname{Cal}}(\widetilde\psi_1)
 &=\overline{\operatorname{Cal}}^{\mathrm{path}}((\lambda^{2^4})_{1/2})
 =\left(\frac12\right)^4
 \overline{\operatorname{Cal}}^{\mathrm{path}}(\lambda^{2^4})\notag\\
 &=\left(\frac12\right)^4 2^4
 \overline{\operatorname{Cal}}^{\mathrm{path}}(\lambda)
 =\overline{\operatorname{Cal}}^{\mathrm{path}}(\lambda)=C_1.
 \label{eq:22-wild-tail-calabi-rescaled}
\end{align}
显然，式~\eqref{eq:22-wild-tail-calabi} 与
\eqref{eq:22-wild-tail-calabi-rescaled} 相矛盾。
\end{proof}

